\documentclass[11pt]{article}
\usepackage[T1]{fontenc}
\usepackage{lmodern,microtype,amsmath,amssymb,amsthm,mathtools,booktabs}
\usepackage[margin=1.04in]{geometry}
\usepackage[unicode,colorlinks=true,linkcolor=blue!55!black,citecolor=blue!55!black,urlcolor=blue!55!black]{hyperref}
\usepackage{xcolor}
\hypersetup{pdftitle={Exact uniform generation of labeled Latin squares},pdfauthor={Mira},pdfsubject={Exact random generation and bipartite perfect matchings}}
\usepackage{enumitem}
\setlist{itemsep=.22em,topsep=.4em}
\setcounter{tocdepth}{1}
\allowdisplaybreaks
\newtheorem{theorem}{Theorem}[section]
\newtheorem{lemma}[theorem]{Lemma}
\newtheorem{proposition}[theorem]{Proposition}
\newtheorem{corollary}[theorem]{Corollary}
\theoremstyle{definition}\newtheorem{definition}[theorem]{Definition}
\newtheorem{remark}[theorem]{Remark}
\DeclareMathOperator{\per}{per}
\DeclareMathOperator{\tr}{tr}
\DeclareMathOperator{\KL}{KL}
\DeclareMathOperator{\Var}{Var}
\DeclareMathOperator{\polylog}{polylog}
\newcommand{\E}{\mathbb E}
\newcommand{\cL}{\mathcal L}
\newcommand{\one}{\mathbf 1}
\newcommand{\ind}{\mathbf 1}
\newcommand{\inj}{\mathrm{inj}}
\newcommand{\RR}{\mathbb R}
\newcommand{\Tt}{^{\mathsf T}}
\newcommand{\Jm}{\mathbf J}
\newcommand{\st}{\colon}
\title{Exact uniform generation of labeled Latin squares}
\author{Mira}
\date{November 19, 2026}
\begin{document}
\maketitle
\begin{abstract}
We construct an algorithm that returns each labeled Latin square of order $n$ with exactly the same probability. For every $n\ge1$ it terminates almost surely and uses $\widetilde O(n^{25})$ expected bit operations. Choosing rows uniformly from the currently legal permutations gives a biased distribution: a square's probability depends on how many choices remain along the way. We use this dependence in reverse. By designating extra matchings early, we give a square more opportunities to be proposed when its final rows would be costly. With three rows left, a cap on the weight of a matching and two early correction blocks control the otherwise exponential variation. A rejection procedure combines the resulting histories at several scales so that each square has the same output probability. The algorithm requires neither a Latin-completion counting oracle nor a permanent oracle.
\end{abstract}
\section{Introduction}
A Latin square has $n$ rows, $n$ columns and $n$ symbols; each symbol occurs exactly once in each row and column. All three sets here are identified with $[n]=\{1,\ldots,n\}$. Let $\cL_n$ be the set of these labeled squares. We do not identify relabelings, transposes or isotopes. Retaining the row labels is essential: our construction randomly orders them to analyze a fixed target, then restores all their labels.

\begin{theorem}\label{thm:main}
There is an algorithm using independent fair bits which, for every integer $n\ge1$, terminates almost surely and outputs each square in $\cL_n$ with probability $1/|\cL_n|$. Its total expected number of bit operations, including preparation, failed trials and restarts, is $\widetilde O(n^{25})$, where a tilde suppresses fixed powers of $\log(2n)$.
\end{theorem}
Time is measured as a function of the order $n$, not its binary encoding length. The theorem gives an expected bit bound, not a worst-case guarantee or a prediction of practical speed. Earlier exact samplers include the probabilistic divide-and-conquer algorithms of DeSalvo \cite{DeSalvo}. Our construction does not complete arbitrary partial squares or establish rapid mixing of a Latin-square walk such as the chain of Jacobson and Matthews \cite{JM}.

Here is the bias the algorithm must remove. Regard a column--symbol pair as an edge of the complete bipartite graph $K_{n,n}$. A row is a perfect matching, so a square is an ordered factorization into $n$ such matchings. For a board $A$ of allowed edges, write $W(A)=\per A$ for its number of perfect matchings. If each new matching is chosen uniformly from the unused edges, a particular ordered square has probability $\prod_dW(A_d)^{-1}$, where $A_d$ is its unused $d$-regular board. Even the two-row residue can cause a large change: its $c$ alternating components give $2^c$ possible next matchings. This is $2$ on one Hamilton cycle, but $2^{n/2}$ on disjoint quadrilaterals when $n$ is even. A worst-case upper envelope at every individual step is therefore a poor way to remove the bias.

To see why an early matching helps, fix a square $S$ and put its rows in a uniformly random order. The union of its last $b$ rows has the same law as the union of any designated $b$ early rows. A large matching count at the end makes $S$ unlikely under ordinary forward generation. But a large matching count in the early block gives us many distinct ways to have designated a matching there first. Their logarithmic effects have opposite signs and cancel in expectation. The blocks need not be independent. We use each block once; the algorithm never has to decide whether $S$ is exceptional.

The proof separates a probability question from a counting question. Given a proposed square, we can draw a new order and new marks on that square alone. This \emph{reverse experiment} does not generate or enumerate other Latin squares. Sections~\ref{sec:coverage} and~\ref{sec:lookahead-cubic} compare its probability of each history with that of generating the same history forward. They bound the mean of the logarithm of this ratio for \emph{each fixed target}. For efficiency we also need a common scale: it must be large enough for that bound, but not much larger than the number of squares. The next section obtains the scale through regular-board enumeration; Appendix~\ref{sec:incidence} gives the companion count lower bound. Section~\ref{gate:multi-v2} turns a mean log bound, rather than a pointwise envelope, into exactly equal probabilities, and Section~\ref{sec:bit-runtime} counts the required bit operations.
\section{Notation and scale}
Let $W(A)=\per A$, $N_d$ be the number of simple $d$-regular boards on the two prescribed shores, $U_d$ their uniform law, and $m_d=\E_{U_d}W(A)$. Write $Q_d$ for the law of a $d$-board reached by deleting successive uniform matchings from $K_{n,n}$. Neither law is presumed to be the law of the board in a uniform Latin square. Put
\[
 D=\lfloor n^{1/200}\rfloor,\qquad \ell=2+\lceil\log_2(2n)\rceil,
 \qquad J=\one\one^\mathsf{T}/n,\qquad V_d=n!(d/n)^n.
\]
Write $\beta_d=d(1-1/d)^{d-1}$ for $d\ge2$, $\beta_1=1$,
and $L_d=d!\beta_d^{n-d}$. The degree-sensitive permanent inequality of Gurvits \cite{Gurvits,LS} gives $W(A)\ge L_d$ for every $d$-regular board. The van der Waerden and Br\'egman bounds give $V_d\le W(A)\le(d!)^{n/d}$.

The proof uses three comparisons, with constants independent of $n$ and of the target:
\begin{align}
 X_n&=n^{-1}(n!)^{2n}n^{-n^2}(2\pi n/e)^{n/2},
 &\prod_{d=1}^n m_d&=\Theta(X_n),\label{eq:scale}\\
 |\cL_n|&\ge c\prod_{d=1}^n m_d,
 &\sup_{S\in\cL_n}\E_{R_S}\log T&\le\log X_n+O(1).\label{eq:countlb}
\end{align}
Here $R_S$ is a normalized law of auxiliary histories for a prescribed square, and $T$ is its probability divided by the forward probability of proposing that decorated square. These means need not be computed by the algorithm. The second comparison, Proposition~\ref{prop:count-lower}, counts surviving paths of the uniform and matching-biased board measures. The last, Proposition~\ref{prop:mean-cover}, is a bound on \emph{every} square; it could not be replaced by an average over the outputs of the forward algorithm. Together they allow a single upward-rounded rational threshold $H=\Theta(X_n)$, for which $H/|\cL_n|=O(1)$ and $\E_{R_S}\log(H/T)\ge0$. They also show $|\cL_n|=\Theta(X_n)$: by Jensen, $H\sum_w Q(S,w)\ge1$ for each $S$, where $Q$ denotes the forward subprobability, and the forward masses sum to at most one. This gives no asymptotic equivalence with a specified multiplicative constant.
\subsection{Local statistics of a decreasing board}
For a fixed simple bipartite graph $H$ with $e$ edges, let
\[
  \alpha_H=\sum_{v\in V(H)}\binom{\deg_H(v)}2.
\]
In an application $H$ has fixed labeled distinct vertices on each shore. The case of isolated vertices is understood by adjoining their arbitrary distinct labels. Write $\inj(H,A)$ for the number of injective shore-preserving embeddings. A healthy $s$-regular board satisfies $\|A/s-J\|_{2\to2}\le\ell^3/\sqrt{s}$.

The analytic expansion in Proposition~\ref{an:expansion}, obtained from the centered rook representation (cf.\ \cite[Sec.~3]{DW}), expresses $\log W(A)-\log V_s$ on healthy boards as a deterministic fixed sum of normalized finite graph counts, a common correction $\tfrac12(1/s-1/n)$, and an error
\begin{equation}\label{eq:residual}
 O\!\left(\ell^{C}/s^2+1/(s\sqrt n)+n^{-5}\right).
\end{equation}
The fixed truncation order may be large. For $s\ge D$, choose it so that $n(C_h\ell^{C})^h/s^h$ is $O(n^{-5})$. All-forest terms are degree determined. For the other terms, after multiplying by $n^{-a}s^{-j}$, the reference magnitude $n^{-a}s^{-j}(n)_{v_L}(n)_{v_R}(s/n)^e$ is bounded by a constant. Here $(u)_k=u(u-1)\cdots(u-k+1)$. The gradient estimate, for any allowed edge $f$ under a uniform matching $M$, is
\begin{equation}\label{eq:gradient}
 \Pr(f\in M\mid A)=s^{-1}+O(\ell^Cs^{-5/2}+n^{-A_0}),
\end{equation}
where the fixed exponent $A_0$ can be chosen as large as necessary. Two disjoint pins occur with probability $s^{-2}(1+O(\ell^C/s))$, and any fixed $k$ disjoint pins with probability at most $C_k/s^k$. These are estimates on an actual perfect matching, not on independent edges.

For a surviving embedding, inclusion--exclusion up to triples gives the probability of surviving the next deletion, divided by $(1-1/s)^e$, as
\[
  1-\alpha_H/s^2+O_H(\ell^C/s^{5/2}+n^{-A_0}).
\]
Only disjoint pairs can disappear together, explaining the coefficient. Iterating yields
\begin{equation}\label{eq:surv}
 \Pr_{Q_s}(H\subseteq A)/(s/n)^e
       =1-\alpha_H(1/s-1/n)+O_H(\ell^C/s^{3/2}).
\end{equation}
An exceptional unhealthy board has probability $O(n^{-A_0})$ by stopped matrix concentration \cite[Cor.~7.5]{Tropp}: in the normalized residual, one deletion contributes a martingale increment of norm at most $2/(s-1)$; the conditional drift from \eqref{eq:gradient} is $O(\ell^C s^{-5/2})$. Up to degree $s$ its variance parameter is $O(1/s)$ and the total drift is $O(\ell^C/s^{3/2})$. Matrix Azuma applied to the self-adjoint dilation, with a union bound over degrees, closes the healthy stopping threshold. With fixed templates $A_0$ can absorb all inverse powers introduced by dividing a pin probability by $(s/n)^e$.

The finite coefficients in this expansion and the common first correction are used in the proof, not as numerical surrogates for the subsequent exact coins.
\subsection{Comparing the high-degree means}
\begin{lemma}\label{lem:high-scalar}
At the fixed cutoff $D$,
\[
 \sum_{d>D}\bigl(\E_{Q_d}\log W-\log m_d\bigr)\le o(1).
\]
This is a comparison of sums of means; no uniform closeness of the decreasing-board law to $U_d$ is required.
\end{lemma}
\begin{proof}
It suffices, by Jensen, to compare the two mean log permanents. Split at $R=\lceil\sqrt n\,\ell^{B_0}\rceil$, with $B_0$ as in the dense-strip estimates of Appendix~\ref{sec:incidence}.

For $D<d<R$ the probability under $U_d$ of a fixed set of $e$ cells $F$ is
\[
 (d/n)^e\{1-\alpha_F/d+O_F(d^{-2}+1/n+(d/n)^2)\}.
\]
For clarity this range starts at $D$, below the auxiliary cutoff in the incidence argument. Here is a direct switching verification. Condition on bounded pins $H$ and consider an additional unforced cell $(x,y)$. Put $N=(d-h_x)(d-h_y)$ for the remaining endpoint degrees, and $L_0=nd-2d^2+2d-1+N$. Let $Z=\sum_{a\ne x,b\ne y}A_{ay}A_{xb}A_{ab}$, $Z_0$ its subsum in which $ay,xb$ are not pins, and $R_H=\sum_{ab\in H,\ a\ne x,b\ne y}(1-A_{ay})(1-A_{xb})$. Equality of insertion and removal counts for switches on $2$-by-$2$ subboards gives, with $p_H=\Pr(A_{xy}=1\mid H)$,
\[
 L_0p_H=N-\E_H Z_0-\E_H[A_{xy}(Z-Z_0)]+\E_H[A_{xy}R_H].
\]
In particular $p_H\le2d/n$, also after bounded additional pinning. With $p=d/n$, this bounds the three sums respectively by $O(n^2p^3+|H|p^2)$, $O_H(np^3+p)$ and $O_H(p)$. Thus $p_H=p(1+O_H(1/d+p))$. In $Z_0$ all but $O_H(n)$ triples avoid pin vertices and by this rough estimate have probability $p^3(1+O_H(1/d+p))$; the others use three new entries except for at most $|H|$ triples using two. Therefore $\E_H Z_0/N=p(1+O_H(1/d+p))+O_H(p/n+n^{-2})$. This cancels the $+p$ in $N/(L_0p)=1-(h_x+h_y)/d+p+O_H(d^{-2}+n^{-1}+p^2)$. Iterating the single-cell estimate proves the formula, as each adjacent cell pair is counted once.

Now substitute these probabilities and \eqref{eq:surv} into (A1) under the two laws. The high-confidence spectral hypothesis for a uniform board is provided by the directed matrix version of Tikhomirov--Youssef \cite[Theorem B]{TY}; use any fixed lower degree power smaller than $1/200$. On either bad event, polynomial bounds on all fixed graph counts and on $\log W$ suffice. In reducing a mixed template, eliminate tree homomorphism counts exactly and expand the remaining collisions into injections. After removing $j$ mean-kernel edges there are at most $j+1$ components. A normalized connected unicyclic injection has bounded reference and its first-order correction is $-\alpha_F/d$ in both laws up to
$O_F(d^{-2}+1/n+p^2+\ell^C d^{-3/2})$. A simple connected component with cycle rank at least two costs $O_F(1/n)$ since its number of edges exceeds its number of vertices. Two or more cyclic components cost $O_F(1/n)$ after their joining mean factors: for injections of their disjoint union use the fixed-pin estimate directly (each unicyclic component has equal numbers of vertices and edges). Additional powers of $1/d$ from repeated edges only improve the bounds. The mode expansion and its universal signed $U/(2n)$ term are identical under the two laws except for these finite-template statistics. The sum of the errors is
$O_h(\ell^C/\sqrt D+R/n+R^3/n^2)=o(1)$; $h$ is fixed.

For $d\ge R$ it suffices to compare the quartic trace. More explicitly put $B=A/d-J$, $U=n/d-1$, $g=1/d-1/n$ and $\tau_4=\tr((B^{\mathsf T}B)^2)$. On the real central Gamma window of (A3), the part of the log rook polynomial depending on the board is $z^4\tau_4/4$ up to $O(n\ell^{C_0}/d^3+n^{-10})$, $z=1/x$. To see the form of this reduction, the determinant in (A4) contributes $z^2 U/2+z^4\tau_4/4$; the cubic paired term is $(2/3)z^3\sum_{ij}B_{ij}^3=(2/3)z^3(n/d^2-3/d+2/n)$, and the single diagonal quartics contribute $-(3/2)z^4ng^2$. Subsequent terms in the connected graph expansion of Appendix~\ref{an:appendix} with surplus at least three are $O(n\ell^{C_0}/d^3)$ by the same rooted summation; the expansions of retained kernels beyond the displayed powers cost the same using $\sigma^2\le\ell^6/d$. The fixed $B_0$ may exceed $C_0$. The tails outside the Gamma window are negligible. In the central integral the response to a unit change of $\tau_4$ is $\frac14\E(x^{-4})=\frac14+O(1/d+1/\sqrt n)$: its mode is $1+O(1/d)$ and its log curvature is between two negative constant multiples of $n$. The derivative error times the bound $\tau_4\le U\sigma^2=O(n\ell^6/d^2)$ is summable $o(1)$ after increasing $B_0$. Thus it is enough to compare $\E\tau_4$.

Only the normalized rooted four-cycle count in $\tau_4$ is not degree determined. Under $Q_d$ its expectation is $(1-1/n)^2[1-4(1/d-1/n)]+O(\ell^C/d^{3/2})$ by \eqref{eq:surv}. Under $U_d$ the row-pair resplitting calculation in the dense-strip proof of Appendix~\ref{sec:incidence} gives $n(d-2)(d-1)^2/[d^3(n-2)]$ with error $O(n/d^3+\sqrt n/d^2)$. The latter expression follows on replacing the centered squared codegree sum there by $w_*$ in $\sum_{i\ne j}c_{ij}(c_{ij}-1)/d^4$ and using $\sum_{i<j}c_{ij}=n\binom d2$. The two rational centers differ by $O(d^{-2}+1/(nd)+n^{-2})$. All remainders sum to $o(1)$, proving the comparison.
\end{proof}

\subsection{An explicit factorial scale}
\begin{lemma}\label{lem:midpoint}
The product of matching means satisfies $\prod_{d=1}^n m_d=\Theta(X_n)$.
\end{lemma}
\begin{proof}
At $d\ge n/2$, uniformly,
$\log m_d=\log V_d+(n/d-1)/2+O(1/n)$.
In fact $U=n/d-1\le1$ and $\|A/d-J\|=O(n^{-1/2})$ except with probability $O(n^{-A_0})$ by \cite[Theorem B]{TY} on the complement when $n-d\ge n^{1/10}$, and the deterministic complement bound $(n-d)/d$ otherwise. Thus $\tau_4=O(1/n)$, the preceding quartic reduction leaves $z^2U/2+O(1/n)$ on the Gamma core, and $X\sim\operatorname{Gamma}(n+1,n)$ has $\E(X-1)=1/n$ and $\E(X-1)^2=O(1/n)$. Taylor's theorem under the positive integral gives $\per A/V_d=e^{U/2}(1+O(1/n))$. On bad boards the Br\'egman bound divided by $V_d$ is $n^{O(1)}$ uniformly in this range, so choose $A_0$ larger than its exponent before taking expectations.

Deleting a marked matching gives $N_dm_d=N_{d-1}m_{n-d+1}$. If $t=\lfloor n/2\rfloor$ and $k=n-t$, complementation gives $N_t=N_k$ and exactly
\[
 \prod_{d=1}^n m_d=
 \frac{\prod_{d=t+1}^n m_d\prod_{d=k+1}^n m_d}{N_t}.
\]
At the median density, the dense equal-margin enumeration of Canfield and McKay \cite{CM} gives $N_t=\Theta(\binom nt^{2n}/\binom{n^2}{nt})$. With $H_j=\sum_{i=1}^j1/i$ the resulting logarithm is
\[
 2n\log(n!)-n^2\log n+\log\binom{n^2}{nt}-n\log\binom nt+
 \tfrac n2(2H_n-H_t-H_k-1)+O(1).
\]
For $\theta=t/n$ the log binomial difference is
$\frac n2\log(2\pi n\theta(1-\theta))-\frac12\log(2\pi n^2\theta(1-\theta))+O(1)$.
Since $2H_n-H_t-H_k=-\log(\theta(1-\theta))+O(1/n)$, the remaining terms combine to $\log X_n+O(1)$ for both parities.
\end{proof}

\section{Proposing histories for a fixed square}\label{sec:coverage}
In this section $S$ is a \emph{fixed} labeled target. Each row is a color on the complete bipartite graph: edge $(i,j)$ has the unique color of the row placing symbol $j$ in column $i$. Write $T_b$ for the union of a uniform $b$-subset of its row colors. A proposal first draws a uniform order of the labeled rows. The reverse experiment on $S$ draws the same kind of order, independently, and then its specified auxiliary marks. Failed or conflicting trials are null outcomes; a step is not conditioned on passing its checks. This distinction determines the normalization.

\subsection{Paying for colored cycles}
We give the details by which the signed finite-pattern bounds become an algorithm. For an injective simple pattern $H$, let $N_H=(n)_{v_L}(n)_{v_R}$ and let $P_j=(d)_j/(n)_j$. An embedding partitions its $e$ edges by equality of their Latin row colors. For a partition $\rho$ let $M_\rho$ be the number of embeddings realizing precisely this partition, $b=|\rho|$, $t=e-b$, and $m_2(\rho)=\sum_{B'\in\rho}\binom{|B'|}{2}$. There are $u=\binom e2-\alpha_H$ disjoint pairs of edges. For a pair of disjoint edges the number of complete-board injections giving them the same color is exactly $N_H/(n-1)$: once their two left endpoints and the first right endpoint are fixed, exactly one of the remaining right endpoints works. Consequently
\begin{align}
\E\inj(H,T_d)={}&N_H\left[P_e+\frac{u}{n-1}(P_{e-1}-P_e)\right]\nonumber\\
 &+\sum_{\rho:t\ge2}
  [P_b-P_e-m_2(\rho)(P_{e-1}-P_e)]M_\rho.\label{eq:partition}
\end{align}
The first line has the same first-order correction $1-\alpha_H(1/d-1/n)$ as \eqref{eq:surv}, with relative remainder $O_H(d^{-2})$. The remaining coefficients may have \emph{either} sign; in absolute value they are at most $(2+e^2)P_b$.

Delete the $a$ mean-kernel edges of a connected mixed template. Suppose the surviving injective simple graph has $c$ components, $e$ edges, and $j\ge e$ original non-mean edge factors. Then $c\le a+1$. Forests are eliminated by their exact homomorphism counts $n^c d^e$, subtracting lower-vertex quotients to enforce injectivity. For a remaining cyclic $H$, select a simple cycle in one cyclic component. Once its colors are specified, an embedding of that component is determined by the image of a root; each other component contributes at most $n$ choices. For uniform distinct names of the $b$ colors, let $X_W$ count the injective rooted embeddings of the selected cycle with colored word $W$. The $t\ge2$ terms of \eqref{eq:partition}, after normalization by $n^{-a}d^{-j}$ and summation over $d>D$, are therefore bounded by
\[
  C_H n^{c-1-a}\sum_{d>D}d^{b-j}\,\E X_W
                    \le C_H D^{-1}\E X_W .
\]
If the selected cycle contains a singleton color, its expected $X_W$ is bounded by $n/(n-b+1)$: each starting image and the other color names determine the one closing color, if it exists. Hence such terms cost only $o(1)$. Otherwise the word alternates between shores, is proper at each site and every color occurs at least twice. Only finitely many words occur at our fixed cutoff order. This argument works for disconnected patterns and does not replace their statistics by a product of independent cycle counts.

\begin{lemma}[Marked polygon head]\label{lem:head}
There exists a polynomial expected-work prefix experiment, of bounded physical length independent of $n$, such that the generic proposal obtained by appending ordinary uniform-matching deletions satisfies on every fixed target
\begin{equation}\label{eq:generic}
 \E_{\rm rev}\log T_{\rm gen}
       \le \sum_{d>D}\log m_d+\sum_{b\le D}\E\log W(T_b)+O(1).
\end{equation}
Here $T_{\rm gen}$ is the ratio of the conditional reverse atom to its forward atom, including the marks. The experiment does not evaluate any of the target-dependent expectations in the display.
\end{lemma}
\begin{proof}
Give every fixed necessary canonical cyclic word, the word $1212$, and the empty choice fixed positive probabilities. Here $1212$ is the four-cycle whose two alternating row colors are repeated. A word of length $2l$ using $b$ colors prescribes the first $b$ row slots, one per color. Write $e_j$ for its occurrences of color $j$. The marked objects consist of $m$ vertex-disjoint \emph{rooted, labeled} copies of this polygon together with raw permutations extending them. The total number is
\[
  Z_W=(n!)^b\sum_{m=0}^{\lfloor n/l\rfloor}
   \frac{(n)_{ml}^2}{m!\prod_{j=1}^b(n)_{m e_j}} .
\]
Choose $m$ by these exact rational weights. Assign distinct vertices to the ordered copies on each shore and finish each row by a uniform permutation on the unassigned vertices; forget the order of copies. Every decorated object before veto has mass $1/Z_W$. Properness gives $e_j\le l$, hence the vector $(ml,ml,0,\ldots)$ majorizes $(m e_j)_j$. Concavity of $k\mapsto\log(n)_k$ gives $Z_W\le e(n!)^b$. Reject conflicting rows. Within the assigned rows compute a deterministic maximal vertex-disjoint packing $\mathcal P_W$ of realized labeled copies; reject decorations not contained in that packing. There are at most $n$ prospective root images for the fixed oriented word. Each packed copy conflicts with at most $(2l)^2$ others, since prescribing the image at any position determines a walk. Thus $X_W\le[(2l)^2+1]|\mathcal P_W|$.

In reverse the marks are an independent fair subset of $\mathcal P_W$. Conditional on the public word, the head ratio is $Z_W/2^{|\mathcal P_W|}$; no tiny word probability is inverted. For all sufficiently large $n$, the fixed prior times its per-copy credit dominates the \emph{sum} of all $O_H(D^{-1})$ charges for that same word. Conditional on the word alone, the reverse marks sum to one, so each ordinary residual outside the prefix is the union of a uniform subset of target rows. The nonnegative charges for a truncated interval of degrees are bounded by those for the full interval. Replacing its missing ordinary top $b$ degrees by $b$ raw factorials costs only $O(b^2)$ by $W(A_{n-i})\ge n!(1-i/n)^n$. Combining these inequalities with \eqref{eq:partition}, the finite expansion and Lemma~\ref{lem:high-scalar} proves \eqref{eq:generic}.
\end{proof}


\section{A single proposal with bounded mean excess}
\label{sec:lookahead-cubic}

Write $W(G)=\per G$, $m_b=\mathbb E_{U_b}W$, where $U_b$ is the
uniform simple $b$-regular bipartite board, and put
$\beta_b=b(1-1/b)^{b-1}$ for $b\ge2$. For a fixed target square $S$, let
$T_b$ be the union of a uniform $b$-subset of its rows and
$\mu_b(S)=\mathbb E\log W(T_b)$. All logarithms are natural.
We use the cutoff $D=\lfloor n^{1/200}\rfloor$ from Section~2. It suffices to
describe large $n$; finitely many smaller orders may use exhaustive sampling.

Use one uniform hidden ordering of physical row labels. Begin with the
bounded polygon head, then a two-row block and a three-row block, both described below.
Next place exactly one marked block of each size $b=4,\ldots,D$, in increasing
order. The total number of previously placed rows is $q=O(D^2)$ throughout
this prefix. Continue with ordinary uniform matchings until just three rows
remain, and use the cubic sampler below. Collisions in a raw block or
infeasible masked choices are null outcomes, \emph{not} reasons to repeat the
block. Write $T$ for the ratio of reverse probability to forward subprobability
of a supported history of a given target. All reverse marks are made after
choosing the hidden row order. Different $b$-sets used in the same history
need not be independent.

\subsection{Lookahead in a marked block}

Choose an auxiliary matching $M$ first. Partition its domain sites among
the $b$ distinct physical colors, with owner counts $(v_1,\ldots,v_b)$
distributed proportionally to $\prod_j(n-v_j)!/v_j!$ on $\sum v_j=n$;
conditional on counts use the uniform assignment. For a raw block, choose
$M$ uniformly from all permutations and fill each physical color's
complement uniformly. Every legal decorated block then has forward mass
$1/Z_b$, with the integer
\begin{equation}\label{eq:v2-Zb}
 Z_b=(n!)^2\sum_{\sum_j v_j=n}\prod_{j=1}^b\frac{(n-v_j)!}{v_j!}.
\end{equation}
Indeed the owner-assignment probability is $\prod_j v_j!/n!$, and the
uniform free fills have product probability $\prod_j 1/(n-v_j)!$.
A convolution over partial sums computes the normalizer and draws the
counts; integerizing the single-color coefficients by $n!$ suffices.
In the reverse experiment, draw $M$ uniformly from the union of the
$b$ real rows. On a simple block, the owner of each edge is unique.

Take $b_0$ a sufficiently large fixed integer. For $4\le b<b_0$ use raw
blocks with a conflict veto. All these slots, including their predecessors,
are bounded in number independently of $n$. At $b\ge b_0$ draw $M$
uniformly in the current host, hence avoiding the $q$ preceding physical
rows. At color $j$, its $v_j$ owned edges are fixed; write $X_j,Y_j$
for the free domain and range, $s_j=n-v_j$, and $R_j=X_j\times Y_j$.
In this box forbid preceding physical edges and the other reserved mark
edges. If $\min_j s_j<n/2$, take uniform matching fills in these masks.
Otherwise assign the rational weight
\[
 \phi_j(F)=\prod_{e\in F}\prod_{\substack{k>j\\ e\in R_k}}(1-1/s_k)
\]
to an allowed matching $F$ of the free box, and draw proportionally.
This is a forward choice depending only on information already drawn.
The reverse experiment is unchanged. Write $P_j$ for the total weight in
the $j$th mask. On a balanced supported history, the incremental log
reciprocal \emph{relative to using ordinary matching draws in these same
slots} is
\begin{equation}\label{eq:v2-inc}
 \log\frac{Z_b}{(n!)^b}-\log W(T_b)
 +\log\frac{W(B_q)}{n!}
 +\sum_{j=1}^b\left[
     \log\frac{P_j}{s_j!}-\log\phi_j(F_j)
                         -\log\frac{W(B_{q+j-1})}{n!}\right].
\end{equation}
Here $B_i$ means the host after $i$ physical rows; the auxiliary matching
does not consume a row. Formula \eqref{eq:v2-inc} follows by replacing
the raw proposal atom for the mark and each free fill by its actual
probability. A dense unweighted slot has inverse probability $W(B_i)$.

\begin{lemma}[Weighted dense first order]\label{lem:v2-weighted-dense}
For an $s$-by-$s$ matrix $0\le w_{ij}\le1$ whose deficit $1-w$ has
maximum row or column sum $K=o(s)$,
\[
 \log(\per(w)/s!)=-\frac1s\sum_{ij}(1-w_{ij})+O(K^2/s).
\]
\end{lemma}
\begin{proof}
Interpolate $W_t=1-t(1-w)$ for $0\le t<1$. In the matching measure of
any such positive board, an edge marginal is at most $1/(s-2K)$:
from a matching containing the edge, swap its column with every other
row. The sum of relative new weights is at least $s-1-2K$ by the
row and column deficit bounds, since the removed weights are at most
one. Every new matching has a unique predecessor. This also bounds
marginals in a one-row/column minor with $s$ replaced by $s-1$.
Let $Z_{iy}$ denote the minor removing row $i$ and column $y$.
Change a matching in that minor into one missing column $z$ by sending
the row presently owning $z$ to $y$. The expected relative weight is
at least $1-K/(s-1-2K)$ by the marginal bound. Interchange $y,z$ for
the converse. Row expansion gives
$Z_{iy}/\per(W_t)=(1+O(K/s))/s$. Integrating the logarithmic derivative
$-\sum(1-w_{iy})Z_{iy}/\per(W_t)$ and taking $t\uparrow1$ proves the
claim, since the total deficit is at most $sK$.
\end{proof}

In $R_j$ the hole set is the disjoint union of exterior rows $Q$, the
$s_j$ cells of $M$ in the box, and all earlier \emph{free} fill edges
$F_i$ inside $R_j$. Earlier \emph{owned} edges are already counted in
$M$. The row and column weight deficits are at most $q+3b+1$.
Expanding the product in $\phi_j$ and applying
Lemma~\ref{lem:v2-weighted-dense} gives
\begin{align*}
 -\log(P_j/s_j!)={}&
 \frac{e_Q(R_j)}{s_j}+1+
 \sum_{i<j}\frac{|F_i\cap R_j|}{s_j}+
 \sum_{k>j}\frac{|R_j\cap R_k|}{s_js_k}
 +O((q+b+1)^2/n).
\end{align*}
Indeed the first-order soft-weight deficit is at most $2b/n$ per cell;
extending its sum to hard holes costs $O(b(q+b+1)/n)$ after division
by $s_j$. The product's quadratic remainder costs $O(b^2/n)$.
Also $-\log\phi_i(F_i)=\sum_{j>i}|F_i\cap R_j|/s_j+O(b/n)$.
Thus \emph{all} the terms recording earlier free-fill edges cancel in
\eqref{eq:v2-inc}. Conditional on the current target block and its
mark, its $q$ exterior colors are uniform among the other $n-b$ colors.
Consequently
\[
 \mathbb E[e_Q(R_j)\mid T_b,M,\text{internal order}]
      =\frac{q}{n-b}\bigl(s_j^2-e_{T_b}(R_j)\bigr).
\]
Since $e_{T_b}(R_j)\le b s_j$ and $\sum s_j=(b-1)n$, summing after
division by $s_j$ gives $q(b-1)+O(qb^2/n)$. This calculation does
not condition on the identities or decorations of earlier blocks.
The regular terms $\log[W(B_i)/n!]$ equal $-i+O(i^2/n)$ by the lemma.
The $bq$ exterior charge cancels exactly to first order.

For $x_i=v_i/n$ the remaining collision charge is
\[
 \Lambda=b+\sum_{i<j}\frac{|R_i\cap R_j|}{s_i s_j}
       = b+\sum_{i<j}\frac{(1-x_i-x_j)^2}{(1-x_i)(1-x_j)}.
\]
Here the owner sets partition each shore. Expanding the summand as
$1-x_i-x_j-x_ix_j+x_i^2x_j^2/[(1-x_i)(1-x_j)]$ proves
\begin{equation}\label{eq:v2-lambda}
 \Lambda\ge \binom b2+\tfrac12+\tfrac12\sum_i x_i^2
          \ge \binom b2+\tfrac12+\frac1{2b}.
\end{equation}
The charge in \eqref{eq:v2-inc}, apart from its first two terms, is
therefore at most $-1/2-1/(2b)+O(b(q+b+1)^2/n)$ in reverse mean. This
remains valid with a negligible additive error when including unbalanced
marks. Indeed, on \emph{any} $b$-regular block the probability of at
least $k=\lceil n/2\rceil$ owned edges of a specified color is at most
$\binom nk u_b^{n-k}/\beta_b^n$, where $u_b=(b!)^{1/b}$: prescribe
$k$ edges, use Br\'egman on the minor and Schrijver on the denominator.
As $u_b/\beta_b=\exp(O(\log(b)/b))$, $u_b\ge b/e$, this probability,
times $b$, is exponentially small uniformly for $b\ge b_0$. On a
reverse-supported unbalanced history every positive integer matching
count and every factorial has polynomial-size logarithm. One applies
the exceptional-event estimate only to the nonnormalizer terms, retaining
the exact $-\mu_b(S)$.

\subsection{One more term in the composition normalizer}

We record the precision needed to sum \eqref{eq:v2-lambda}. For every
fixed $b\ge3$ one has $Z_b/(n!)^b=O_b(\sqrt n\,\beta_b^n)$.
Indeed with $p_j=v_j/n$ the exponential rate in \eqref{eq:v2-Zb}
is $\sum_j[(1-p_j)\log(1-p_j)-p_j\log p_j]$, uniquely maximized
at the uniform point. Pairwise smoothing is strict for unequal
coordinates whose sum is less than one; faces supported on at most two
coordinates have rate zero, less than $\log\beta_b$. The constrained
Hessian at the maximum is negative definite. Stirling gives the prefactor
$O_b(n^{(2-b)/2})$ there; Gaussian summation over $b-1$ coordinates gives
$O_b(n^{(b-1)/2})$. Elsewhere the polynomial prefactor is dominated by
exponential loss. This suffices for the finitely many raw blocks.

Uniformly for $b_0\le b\le D$ the sharper estimate is
\begin{equation}\label{eq:v2-sharp-Z}
 \log[Z_b/(n!)^b]= n\log\beta_b+\tfrac12\log(2\pi n)
 +O\left(b^{-3/2}+\frac{b\ell^2}{n}+
                 \frac{\ell^3}{\sqrt{bn}}\right).
\end{equation}
To see this, let $(N_j)$ be multinomial with $n$ draws and equal cell
probabilities $1/b$, $\mu=n/b$, $U=\sum(N_j-\mu)^2/n$ and
$\lambda=b/[2(b-1)]$. The exact identity is
\[
 Z_b/(n!)^b=n!(b/n)^n\,\mathbb E\prod_j n^{N_j}/(n)_{N_j},
\]
where $(n)_t=n!/(n-t)!$. On $\max|N_j-\mu|\le\sqrt\mu\,\ell$, expand
$F(x)=x+(1-x)\log(1-x)$ about $1/b$ in the Stirling formula for this
product. Its log equals
\[
 bnF(1/b)+(b/2)\log(1-1/b)+\lambda U
             +O\bigl(b\ell^2/n+\ell^3/\sqrt{bn}\bigr).
\]
Here both linear sums vanish, $F''(1/b)=b/(b-1)$, and $F'''$ is bounded
on $[0,1/2]$. To control the complement, for every fixed $t>0$,
$\mathbb E e^{tU}=O_t(1)$ if $b_0$ is sufficiently large. Split positive
and negative deviations by Cauchy--Schwarz; negative association of
multinomial counts bounds the increasing products and the decreasing
products separately by the products of their one-dimensional expectations.
For either sign of a binomial deviation $Z=N_j-\mu$, integration of the
Chernoff bound $\Pr(|Z|\ge z)\le 2\exp[-z^2/((2+B)\mu)]$ for
$z\le B\mu$ gives $\mathbb E e^{2t(Z_\pm)^2/n}\le1+O_t(1/b)$ on
this range when $b$ is large. For $z>B\mu$, take $B$ large so the
upper tail is at most $\exp[-(2t+2)z]$; the negative tail stops at
$\mu$, and $z^2/n\le z$. The exceptional contribution is negligible
since $\mu\gg\ell^2$. For the original factorial product with all
$N_j\le n/2$, use its Riemann upper bound $\sum nF(N_j/n)\le
bnF(1/b)+U$, then Cauchy--Schwarz and the central-band tail
$O(b e^{-c\ell^2})$. If any $N_j>n/2$, use $\sum F(N_j/n)\le1$
and probability at most $b(4/b)^{n/2}$. These bounds justify the same
expansion after averaging.

For completeness $\mathbb EU=1-1/b$ and
$\Var(U)=2(b-1)(1-1/n)/b^2$: write
\[
 U-\mathbb EU=\frac2n\sum_{i<j}(\mathbf1_{\{B_i=B_j\}}-1/b)
\]
with independent uniform letters $B_i$. Different pairs are uncorrelated.
The sixth centered moment is $O(b^{-3})$. In the six-edge expansion a
draw-vertex appearing just once kills the expectation. For $v\le6$
vertices and $c$ connected components, expansion of the remaining centered
indicators bounds the absolute expectation of each product by
$C b^{-(v-c)}$: the rank of selected equality constraints plus the number
of omitted edges is at least the full rank. Multiplication by
$O(n^{v-6})$, using $v-c\ge v/2$, gives the claim. With the exponential
moment bound, Cauchy--Schwarz bounds the third Taylor remainder. Thus
\[
 \log\mathbb E e^{\lambda U}
    =\tfrac12+\frac{1-1/n}{4(b-1)}+O(b^{-3/2}).
\]
The $1/b$ term cancels the one in
$(b/2)\log(1-1/b)=-1/2-1/(4b)+O(b^{-2})$. This proves
\eqref{eq:v2-sharp-Z}, including absence of an unspecified $O(1/b)$.

The sparse one-planted-matching estimate gives, uniformly in this range,
\[
 \log m_b=n\log\beta_b+\tfrac12\log(2\pi n)-\tfrac12
                 +\tfrac12\log(1-1/b)+O(b^3/n+1/n).
\]
Combining it with \eqref{eq:v2-lambda} and \eqref{eq:v2-sharp-Z}, the
increment at grade $b$ is at most $\log m_b-\mu_b(S)+e_b$, where
$\sum_{b=4}^D e_b=O(1)$: the terms $-1/(2b)-\tfrac12\log(1-1/b)$
are $O(b^{-2})$, $\sum O(b^{-3/2})=O(1)$, and the finite-size errors
sum to $o(1)$ since $q=O(D^2)$. At the finitely many initial raw grades
their normalizer bound, constant $q$, and missing-slot comparison give
$e_b=O_{b_0}(1)$. No independence of a front block and a suffix residual
is used in cancelling the two occurrences of $\mu_b(S)$.

\subsection{The last three rows and a double cubic head}

On a cubic board $G$ the complement of a first matching $M$ consists of
$K(M)$ even cycles; it has $C(M)=2^{K(M)}$ ordered two-row completions.
Put $A=n$, $w(M)=\min(C(M),A)$ and $J(G)=\sum_M w(M)$. Select a
uniform matching and accept with probability $w(M)/A$, repeating
\emph{this matching draw} until accepted. Color the complementary cycles
by independent fair orientations. The probability of a specified ordered
factorization $z$ beginning with $M$ is exactly
\begin{equation}\label{eq:v2-cap-cubic}
 p_G(z)=\frac{1}{J(G)\max(1,C(M)/A)}.
\end{equation}
The expected number of uniform-matching calls is at most $A$. Use this
procedure for the final three rows.

For the additional \emph{pair head}, draw a first uniform permutation and
a relative derangement $\pi$ with weight
$\max(1,2^{\#\mathrm{cycles}(\pi)}/A)$. Its normalizer divided by
$(n!)^2$ is at most $1+(n+1)/A=O(1)$ since
$\sum_{\pi\in S_n}2^{\#\mathrm{cycles}(\pi)}=(n+1)!$.
The integer cycle recurrence by the cycle of the least remaining point,
stratified by the number of cycles, implements the draw with rational
probabilities. Veto all earlier conflicts. The positive overflow log in
\eqref{eq:v2-cap-cubic} for the actual last two rows cancels in reverse
mean with the negative overflow log from this head, since the two
ordered row pairs have the same marginal law.

The additional \emph{cubic head} draws two ordered 3-factorizations of
the same raw edge multiset. Each of the $n$ vertices on each shore has a
permutation $\sigma\in S_3$ identifying its three physical colors with
three auxiliary colors. Require identical left/right counts
$h_{ij}$ of occurrences of $\sigma(i)=j$. Match the two sets of
$h_{ij}$ vertices bijectively and independently for every $(i,j)$.
Veto multiple edges and collisions with previous rows. Every simple
legal pair has exactly one representation. The unrestricted count is
\begin{equation}\label{eq:v2-cubic-Z}
 Z_3^\star=\sum_h d_n(h)^2\prod_{ij}h_{ij}!,\qquad
 d_t(h)=[x^h](\per x)^t,\quad x\text{ a }3\text{-by-}3\text{ matrix}.
\end{equation}
Uniform raw generation uses the displayed weights and two independent
backward samplers from $d_t(h)=\sum_{\sigma\in S_3}d_{t-1}(h-P_\sigma)$.
Only the upper-left $2$-by-$2$ counts need be stored since margins equal
$t$. This is fixed-dimensional integer arithmetic. On a given physical
cubic block $G$, draw its reverse auxiliary coloring by
\eqref{eq:v2-cap-cubic}. Its density is at most $1/J(G)$, so the
head log reciprocal pays $-\mathbb E\log J(G)$, precisely the mean
positive cost of the terminal cubic sampler.

\begin{lemma}[The raw cubic count]\label{lem:v2-cubic-moment}
 $Z_3^\star/(n!)^3=O(n(4/3)^n)=O(m_1m_2m_3)$.
\end{lemma}
\begin{proof}
For a $3$-by-$3$ matrix $B$ of independent standard complex Gaussians,
orthogonality of monomials gives $\mathbb E|\per B|^{2n}=Z_3^\star$
by \eqref{eq:v2-cubic-Z}. Independent row-radial integration contributes
$[(n+2)!/2]^3$. For three independent unit complex rows $z,u,v$, put
$P_i=u_jv_k+u_kv_j$, $\{i,j,k\}=\{1,2,3\}$. Conditional averaging over
$z$ gives exactly $2\|P\|^{2n}/[(n+1)(n+2)]$ by the
$\operatorname{Beta}(1,2)$ squared-coordinate law. Put $t=\langle u,v\rangle$.
Expansion gives
\[
 \frac43-\|P\|^2=\frac{1-|t|^2}{3}
                 +2\sum_i|\overline{u_i}v_i-t/3|^2
 \ge c\left[1-|t|^2+\sum_i(|u_i|^2-1/3)^2\right].
\]
For the inequality rotate $v$ so $t\ge0$; the two vectors of deviations
on the two sides differ in norm by at most
$\|u-v\|+(1-t)/\sqrt3=O(\sqrt{1-t^2})$.
Conditional on $u$, $|t|^2\sim\operatorname{Beta}(1,2)$; the integral of
$\exp[-cn(1-|t|^2)]$ is $O(n^{-2})$. Squared coordinates of $u$ are
uniform on the two-dimensional simplex, so the integral of
$\exp[-cn\sum_i(|u_i|^2-1/3)^2]$ is $O(n^{-1})$. The total angular
factor, including the first $z$-integral, is $O((4/3)^n n^{-5})$.
Multiplying the radial factor proves the first assertion. For the second,
$m_1=1$, $m_2=\Theta(\sqrt n)$ by degree-two cycle enumeration, and
$m_3=\Theta(\sqrt n(4/3)^n)$ by the sparse one-plant estimate.
\end{proof}

\subsection{Mean coverage, common scale and factor cost}

\begin{proposition}[Uniform mean coverage]\label{prop:mean-cover}
For the single decorated proposal just constructed, every fixed labeled
target has $\mathbb E_{\rm rev}\log T\le\sum_{d=1}^n\log m_d+O(1)$,
with an absolute additive constant (after fixing the degree exponent
and analytic truncation).
\end{proposition}
\begin{proof}
By Lemma~\ref{lem:head}, the total high-degree contribution, including the marked polygon, is at most
 $\sum_{d>D}\log m_d+O(1)$, using the common
first-order saddle correction before summing the errors; an $O(1/d)$
unspecified remainder would not suffice. Compare the new first five
fixed head slots against missing ordinary dense slots at a cost $O(1)$:
each $W(B_i)\ge n!(1-i/n)^n$ with bounded $i$. For each $4\le b\le D$
the ordinary suffix has mean log contribution $\mu_b(S)$ and the new
head block's mean comparison to the same physical slots is at most
$\log m_b-\mu_b(S)+e_b$. Sum using the bound above. For $b=1,2,3$
use the paired cubic and overflow cancellations and
Lemma~\ref{lem:v2-cubic-moment}. These are sums of expectations under
a single normalized reverse law; no favorable intersection of events
is required.
\end{proof}

In particular a sufficiently large fixed multiple of any computable
scale $X_n\asymp\prod_d m_d$ suffices for the soft cap. The
constant-loss incidence lower gives $H/|\mathcal L_n|=O(1)$ for such
a choice. This does not assert a constant pointwise reverse gate
probability. There is a counting consequence independent of the cost
of the gate: the forward marginal on $S$ is
$Q(S)=\mathbb E_{\rm rev}(1/T)\ge\exp(-\mathbb E_{\rm rev}\log T)$.
Since $\sum_S Q(S)\le1$ and the proposal is the same for every $S$,
one obtains $|\mathcal L_n|=O(\prod_d m_d)$. Together the bounds give
$|\mathcal L_n|=\Theta(X_n)$, not an asymptotic equivalent.

Finally, the weighted fill does not require a larger matching graph.
For each future $k$ let $l_k=\max(0,|X_j\cap X_k|+|Y_j\cap Y_k|-s_j)$.
If the chosen matching has $m_k'$ cells in $R_k$, then
$0\le m_k'-l_k\le v_k$. Thus
$\psi(F)=\phi_j(F)/\prod_{k>j}(1-1/s_k)^{l_k}$ lies between $e^{-4}$
and $1$ on balanced boxes: use $\sum v_k=n$ and $s_k\ge n/2$.
Rejecting a uniform matching with this acceptance rule gives the
weighted draw in constant expected attempts. Its normalizer is an
ordinary unweighted matching count times a known rational times the
mean of a $\psi$-coin of bias at least $e^{-4}$. There is one such
scalar per physical fill, $O(D^2)$ in total. For each of the two cubic
boards the only new scalar is the bias
$\theta=J(G)/(A\,W(G))=\mathbb E_M[w(M)/A]\ge1/n$.
These positive factors and their reciprocals use the same
geometric/negative-binomial gate as the matching marginals, without
anchor generation for the scalar itself. No unknown $J$ or weighted
permanent is evaluated to decide a coin.

\section{Exact flattening at several scales}\label{gate:multi-v2}

We work first on an arbitrary finite set of $N$ targets. A trial proposes a
target $s$ and decoration $w$, or returns null. Let $Q(s,w)$ be its
forward atom, and $R_s(w)$ the atom in a normalized reverse experiment
at $s$. Both include exactly the stated random labels, marks, and
sampling choices of the construction; internal runtimes are not part
of the decoration. We discard forward atoms outside reverse support
and require $R_s(w)>0\Rightarrow Q(s,w)>0$. Write
\[
 T(s,w)=\frac{R_s(w)}{Q(s,w)},\qquad z(s,w)=\frac{H}{T(s,w)} .
\]
The deterministic positive rational $H$ is common to all targets.
The previous estimates supply two quite different inequalities:
\begin{equation}\label{gate:hyp-v2}
 \mathbb E_{R_s}\log z\ge0\quad(\text{every }s),\qquad H/N=O(1).
\end{equation}
The procedure below does not evaluate either expectation or know $N$.
It also works with a larger $H/N$. Fix a known integer $M\ge2$
with $\log z\le M$ on every reverse history; no corresponding
pointwise lower bound on $\log z$ is needed. For our proposal a
polynomial $M$ suffices. For example, conditional on the public word
and order (whose probabilities cancel), the reverse experiment takes
$O(n)$ matching or coloring marks, each with atom at least
$1/[(n!)^{O(1)} 2^{O(n)}]$, and $O(n)$ independent binary packing
choices. A capped cubic coloring atom satisfies this same crude
bound. Eligible forward probabilities are at most one, and
$\log H=O(n^2\log n)$. We can use $M=n^4$ for all sufficiently large
$n$, with no unknown coefficient needed in this choice.

\begin{lemma}[Multiscale exact flattening]\label{gate:multiscale-v2}
Suppose a forward trial followed by its soft coin of bias $T/(T+a)$,
and a fresh reverse trial at a prescribed $s$ followed by its soft
coin of bias $a/(a+T)$, each cost at most $W$ in expectation,
uniformly for $a=H/2^j$, $0\le j\le\lceil\log_2M\rceil$.
Assume the mean-log inequality in \eqref{gate:hyp-v2}. There is an
exactly uniform sampler of expected total cost
\[
 O\!\left((W+\log^2(2M))(1+H/N)\log^2(2M)\right)
\]
apart from preparing the common inputs. In particular a polynomial
bound on $M$ does not cost a power of $n$.
\end{lemma}

\begin{proof}
Put $J=\lceil\log_2M\rceil$, $m_j=2^j$, $m=m_J$, $u=1/64$, and
\[
 f_j(s)=\mathbb E_{R_s}\frac{z}{1+z/m_j},\qquad
 g_j(s)=f_j(s)/m_j,\qquad r_j(s)=\frac{u}{u+f_j(s)} .
\]
On scale $H/m_j$ the soft forward proposal gives exactly $f_j(s)/H$
mass at $s$; this is the atom identity
$Q T/(T+a)=R_s/(T+a)$. One fresh soft reverse experiment is a
$g_j(s)$-coin. Define $q_s=\mathbb E_{R_s}(1/T)$ only for the proof.
Then $\sum_s q_s\le1$, $\mathbb E_{R_s}z=Hq_s$, and
$0<f_j(s)\le\min(m_j,Hq_s)$. Jensen also gives $Hq_s\ge1$.

For $v\ge2$ and $t\le v$ we have
\[
 \frac{e^t}{1+e^t/v}\ge \frac{1+t}{3}.
\]
When $t\le\log v$ use $e^t/2$ and $e^t\ge1+t$ (a negative right
side is automatic); otherwise use $v/2\ge(1+t)/3$. Thus
$f_J(s)\ge1/3$ for every $s$.

Choose $j$ uniformly in $\{0,\ldots,J\}$ and run its soft forward
proposal. Upon receiving an $s$, require fresh independent $r_i(s)$
tests to pass for $i=0,\ldots,j-1$, stopping at the first failure.
For $j<J$ require one more $r_j$ test. For $j=J$ instead require a
test of bias $u/f_J(s)$. Output $s$ iff all required tests passed.

To implement each $r_i$, repeatedly pick a fair side, try its
$u/m_i$ or $g_i$ coin, and return which side first succeeded. This
costs $O(m_i)$ coin calls uniformly. At the final scale set $c=u/m$
and use the same race to supply $p=c/(c+g_J)$; $p\le3u/(1+3u)<1/5$.
Draw $k\ge1$ with probability $2^{-k}$ and accept iff $k$ independent
$2p$-coins pass. The resulting bias is
$\sum_{k\ge1}p^k=p/(1-p)=u/f_J$. The bounded linear factory of
Huber \cite{Huber} with this fixed slack needs bounded expected input
calls; it uses only fair bits and rational flips. To see bit overhead
can also be bounded at fixed slack, use the doubling-stage version:
at stage $j$ the barrier is $32\,2^j$, relative slack $2^{-j-1}$,
and the amplifier stays between $2$ and $2e^{1/2}$. The exit walk has
mean drift bounded away from zero for $p\le1/5$ and expected length
$O(2^j)$ (the geometric overshoot has bounded mean); its survival
factor to the next stage is at most
$(1+2^{-j-2})^{-32\,2^j}\le e^{-6}$. Rational parameters at this
stage have $O((j+1)^2)$ bits. Implement barrier powers by sequential
rational coins; the geometrically summable stage costs bound auxiliary
bit work. Thus a final test, if invoked, costs $O(m(W+\log^2(2M)))$.

Before the uniform level prior, the output mass at level $j<J$ is
\[
 (f_j/H)\prod_{i\le j}r_i
     =(u/H)(1-r_j)\prod_{i<j}r_i.
\]
At level $J$ it is $(u/H)\prod_{i<J}r_i$. Summing telescopes:
each target receives $u/[H(J+1)]$, exactly.

It remains to charge expensive tests, not simply multiply by their
worst-case cost. Let $B=\{s:Hq_s>m^4\}$, a proof-only set of at most
$H/m^4$ targets. At any level the first-gate mass of $B$ is at most
$(m/H)|B|\le m^{-3}$. Even running every test thereafter costs
$O(m(W+\log^2(2M)))$ on these targets. On all other $s$ put
$K=\log(8m^4)$. Since $(t-K)_+\le e^{t-K}$,
\[
 \mathbb E_{R_s}\min(\log z,K)\ge -e^{-K}Hq_s\ge-1/8 .
\]
Monotonicity and the scalar inequality applied at
$\min(\log z,K)$ give $f_i(s)\ge7/24>1/4$ when $m_i\ge\max(2,K)$.
Then $r_i\le1/17$. Up to the first such $i_0$ the sum of the
worst-case test costs is $O((W+\log^2(2M))\log(2m))$.
Thereafter reaching $i$ costs a probability factor at most
$17^{-(i-i_0)}$; its conditional work grows only as $2^i$.
This bound includes the final inversion. If there is no $i_0$ then
$m<K$ and the same cost bound holds. The first-gate masses form a
subprobability, so these conditional estimates prove the same bound
per complete trial. Independent restarts have a.s. finite stopping
time; the mean number is $H(J+1)/(uN)$. Renewal gives the stated
total expectation even when cost correlates with the success of an
individual trial.
\end{proof}

\subsection{The soft coin, with all matching probes charged}

\begin{lemma}[Binary-prefix permanent gate]\label{gate:binary-v2}
For the decorated row, weighted-fill and cubic proposals in this
paper, both soft coins in Lemma~\ref{gate:multiscale-v2} can be
implemented in expected bit work
$\widetilde O(n^4P(n)+n^{12})$. Here $P(n)$ bounds the expected
bit cost of a fresh uniform perfect matching on every feasible
\emph{unweighted} bipartite board of side at most $n$. This generous
ancillary allowance need not be optimized. No Latin-square completion
oracle, exact permanent evaluation on the ordinary path, or
exponential-time path enumeration is assumed.
\end{lemma}

\begin{proof}
After null histories have been discarded, $T$ is a known positive
rational multiplier times products and reciprocals of ordinary
permanents and positive coin biases. For a balanced weighted fill,
divide the entry-weight product by its explicit common upper
bound from the free-box intersections; simultaneous attainment of that
bound is not required. The resulting weight of a
uniform permissible matching is a rational in $[e^{-4},1]$.
Rejecting by this weight samples the weighted matching, and its
normalizer is the unweighted permanent times a known factor times
the mean of this coin. The occupancy-range argument that gives this
uniform lower bound is important: repeated sampling is not assumed
cheap merely because the individual entries are positive. Imbalanced
compositions use the ordinary fill. A capped cubic block has the
additional factor $J(T_3)=A\per(T_3)\theta$, where $A=n$ suffices and
$\theta=\mathbb E_M\min(1,2^{K(M)}/A)\ge1/A$ for uniform $M$ on
the cubic board. The relevant $K(M)$ is found by following degree-two
cycles. There are only a bounded number of these latter coin factors.
Neither the (possibly very large) individual number of cycles nor a
sum over all perfect matchings is computed to obtain $\theta$.

For each ordinary permanent factor draw an independent uniform
matching called its anchor. Traverse the left vertices in a fixed
order. At a vertex split the list of possible column indices by the
bits of their index in a fixed, padded binary encoding. Reveal the
anchor's next branch bit. Its conditional probability $p$ can be
probed by taking a fresh uniform matching on the current minor with
that row restricted to the branches already revealed. After the
column is fixed, delete that row and column. The product of the
encountered $p$'s equals $1/\per B$ for \emph{each} anchor on board
$B$. At any visited site, averaging over the anchor's next choice
before observing it gives $\mathbb E(1/p)\le2$, the number of
nonempty children. This statement is not claimed conditional on a
complete rare anchor. Use fixed binary depths, even when one child is
forced. Over $O(n)$ boards and the scalar coins above, there are
$D_0=O(n^2\log n)$ sites. Frequencies from a single matching are not
reused as independent probes.

Fix $k=D_0+2$, or use a deterministic upper bound on $D_0$.
At every site write $T_j'$ for the number of independent $p$-coin
flips to obtain $j$ successes. For a required factor $1/p$ take
$T_k'/k$; for a required $p$ take $(k-1)/(T_k'-1)$.
Their means are $1/p$ and $p$, respectively. Their size-biased
versions are $(T_{k+1}'-1)/k$ and $(k-1)/T_{k-1}'$.
These four identities follow directly from
$\Pr(T_j'=t)=\binom{t-1}{j-1}p^j(1-p)^{t-j}$, including $p=1$.
Their relative second moments before size bias are at most
$1+1/k$ and $(k-1)/(k-2)$, the latter by replacing
$(T_k'-1)^{-2}$ by $[(T_k'-1)(T_k'-2)]^{-1}$. Taking independent
factors times the known multiplier yields $Y>0$ with
$\mathbb E Y=T$ and $\mathbb E Y^2\le eT^2$. Both statements hold
conditional on all anchors. The size-biased law $Y^*$ therefore
retains the original anchor law: replace each local factor by its
displayed biased version, independently. Purely deterministic
ratios can of course be sampled without anchors.

Let $l=\mathbb E\min(1,Y/a)$ and $h=\mathbb E\min(1,a/Y^*)$.
Then $Th=al$. The second-moment bound gives
$l\ge(8e)^{-1}\min(1,T/a)$ by Paley--Zygmund. Thus $l+h\ge(8e)^{-1}$
and an independent fair-side race between $l$ and $h$ has bounded
expected rounds. The winning side gives, respectively, $T/(T+a)$
or $a/(T+a)$. This exactness uses neither a logarithmic estimate of
$T$ nor exact knowledge of any of its unknown factors.

A full race round needs $O(kD_0)$ ordinary matching probes in
expectation by the binary-branch identity. The bounded number of
cubic scalar sites add $O(kn)$ and the bounded-away-from-zero fill
sites are covered by $O(kD_0)$. To obtain the decorated forward or
reverse history itself uses only $O(n)$ ordinary draws and at most
$O(n)$ expected draws for the capped cubic rejection, plus its
constant-factor weighted-fill rejections. An independent round may
be long and may win, but renewal with an unconditional constant
success bound still gives $\widetilde O(n^4)$ probes in total.

Here are generous sequential bit allowances for auxiliary work.
A chosen polygon word has fixed length; enumerate its $n$ rooted
candidates, not all assignments of its vertex labels. The composition
normalizer at a $b$-block is a convolution table with $O(bn^2)$
integer operations on $\widetilde O(bn)$-bit entries after factorial
clearing; all block lengths sum to at most $n$. The cubic
double-color histogram has four free counts of range $n$ at each of
$n$ levels. Its $O(n^5)$ transitions can use streaming scans and
fixed-offset shifts of tables in lexicographic order; a constant
number of sorted streams suffices for the six possible vertex types.
Table entries and terminal factorial weights have
$\widetilde O(n)$ bits, and sequential backward sampling may rescan a
full table at each level. The degree-two overflow pair uses the
ordinary integer cycle-index table. These arithmetic costs fit
$\widetilde O(n^9)$ even with schoolbook products and comparisons.

Explicit normalized weighted-fill coins use at most $O(b)$ rational
factor powers, with exponents at most $n$, so their numerator and
denominator have $O(bn\log n)$ bits. They can be computed by repeated
squaring and schoolbook multiplication in $\widetilde O(n^5)$ bit
operations per probe, even using $b\le n$ and scanning every filled
cell for each future free box. The full known path prefactor has
$\widetilde O(n^3)$ bits by summing over the $O(n)$ physical slots;
including $H$ and its dyadic rescalings does not change this bound.
Feasibility, copying a board, and validating a probe cost a polynomial
scan smaller than these allowances. All probe probabilities at
anchor sites are at least $1/n!$. Uniformly in their realized
values, the uncapped negative-binomial counters have
$\mathbb E\log^2(1+T_j')=O((n\log n+\log k)^2)$ by their geometric
tails or Markov's inequality at exponentially increasing multiples
of their mean. For increment costs one may use size bias and
$\mathbb E[T_j'\log(1+T_j')\mid p]\le
(j/p)\log(1+(j+1)/p)$. Thus the unreduced product's squared bitlength
has expectation $\widetilde O(n^6)$, and even $D_0$ schoolbook
multiplications at that bound are harmless here. All rational
choices can be driven by fair bits by sampling a uniform integer
below their denominator. These bounds prove the displayed allowance;
they are bounds on expected total work, not RAM arithmetic counts.
\end{proof}

With $M$ polynomial, $H/N=O(1)$ and the conservative matching
adapter, the preceding two lemmas give the theorem's total bit
bound after adding the one-time scale preparation. An explicit
constant times the calibrated scale may be replaced by an extra
factor $\log(2n)$ for all sufficiently large orders if its implicit
absolute constant is not numerically evaluated; this changes only
logarithmic factors in the bound. The large-order inequalities need
not hold on an unspecified small input. Fix a sufficiently large
admission threshold and use exact exhaustive uniform generation
below it, including $n=1$. This finite modification does not affect
the exponent. It is not a claim that a tested, practical numerical
admission cutoff or a scalable implementation is supplied.

\section{Expected bit time}\label{sec:bit-runtime}
Here $\widetilde O$ suppresses fixed powers of $\log(n+2)$.  We count
individual unbiased random bits and bit operations, not exact arithmetic on
unbounded integers at unit cost.  It suffices to use an unweighted matching
sampler.  All stated exponents are for the order $n$, not its encoding length.

\begin{lemma}[A conservative finite-bit matching clock]\label{lem:pmclock-v2}
For a feasible zero-one board of side $s\le n$, an exact uniform
perfect matching costs $P(n)=\widetilde O(n^{21})$ expected bit operations.
\end{lemma}
\begin{proof}
We specify the bit-accounting adapter of the finite-table correction in
\cite[Sec.~6]{GLMP}.  The clock below is deliberately conservative.  Use
the lazy single-edge Jerrum--Sinclair--Vigoda chain on perfect and
near-perfect matchings of the positive completion of the board.
From perfect, propose removing a uniform matched edge; from near-perfect,
choose a uniform vertex for the corresponding plug or swap.  Use one-half
laziness and Metropolis acceptance.  Factor-two accurate hole weights
give inverse gap at most $1000s^6$ for $s\ge2$: the cross-flow bound
is $48s^4$, and the full-flow multiplier is at most $2+4(r+1/r)$ for
$r\in[s^2/2,2s^2]$, the near-to-perfect stationary mass ratio
(\href{https://www.math.cmu.edu/~af1p/Teaching/Markov_Chain_Mixing/Papers/permanent.pdf}{Jerrum--Sinclair--Vigoda,
Lemmas 4.1, 4.4}). Start each independent chain at a deterministically
found real matching $x_0$ of stationary mass at least
$[(1+2s^2)s!]^{-1}$ on good weights; side zero or one is immediate.

Let $u=C(s+2)^4(1+\lceil\log_2(s+2)\rceil)^2$ rounded up,
$\delta=2^{-u}$, with a sufficiently large fixed $C$.
At $R=O(s^2\log(s+2))$ temperatures, cool the fake edges incident to one
left vertex by the rational factor $3/4$, stopping exactly at $1/s!$.
A configuration has at most one affected edge, so adjacent ideal hole
ratios differ by at most $4/3$. At each update draw independently
$S=O(s^2(u+\log(s+2)))$ end states of
$T_0=O(s^6(s\log(s+2)+\log(s+2)))$ steps. Update each hole weight
by its old value times the ratio of the observed full to the corresponding
near-class frequency; a missing denominator or full class fails the attempt.
The initial weights are $s$.  With sufficiently large absolute coefficients,
the reversible $L^2$ mixing bound from $x_0$ gives TV error less than
$1/(1000s^2)$ per preliminary end state on good weights.  Each class
has stationary probability at least $1/[2(1+2s^2)]$ there. Independent
binomial concentration at fixed relative accuracy (say $1/50$) and a
union bound over all $R(1+s^2)$ classes give failure at most $2^{-20u}$
for the complete calibration. Indeed the ideal update is exactly
$w\pi(\mathrm{full})/\pi(\mathrm{hole})$; empirical accuracy within
a multiplicative factor $6/5$, followed by cooling, maintains the
factor-two condition.  This argument is inductive and does not multiply
per-temperature estimation errors.

At the end of this attempt draw one fresh specified chain of
$T_1=O(s^6(s\log(s+2)+u))$ steps.  On good weights the TV
error is at most $2^{-20u}$, and conditional on a real perfect endpoint
the \emph{stationary} law is exactly uniform: real matchings all have
weight one. The stationary probability of real endpoints is at least
$1/[2(1+2s^2)]$, since all fake perfect matchings together have weight
at most $s!/(s!)=1$. Repeat \emph{independent entire attempts} until
success, at most $J=O(s^2u)$ times, returning $x_0$ if exhausted.
Success has probability $\Omega(s^{-2})$, hence the expected number of
attempts is $O(s^2)$, not $J$. Bad calibration and finite mixing
together perturb the conditional success law by at most $O(s^2 2^{-20u})$.

Round token probabilities by dividing a uniform $b$-bit integer into
$4s$ almost equal bins, and round each binary acceptance probability down
to a multiple of $2^{-b}$.  Take $b=64u+O(\log(2+J(RST_0+T_1)s))$ with
sufficient padding.  Coupling all draws across the \emph{whole capped}
horizon costs at most $2^{-60u}$ in TV.  Adjusting absolute constants, the
resulting finite-output law $Q$ has TV error at most $2^{-8u}$ from the
uniform law $\pi$. Since every valid atom has $\pi(x)\ge1/s!$,
$\pi-(1-\delta)Q$ is nonnegative.  With probability $1-\delta$ use
the capped sampler; otherwise compute and sample the correction
$[\pi-(1-\delta)Q]/\delta$.

For completeness this last law is attainable at negligible \emph{expected}
cost. A single hole weight is a product of at most $R$ frequency ratios
with integers at most $S$. There are at most $(S+1)^{2Rs^2}$ weight
vectors. For each phase and vector, compute one end-state distribution
by time-layer dynamic programming on at most $(1+s^2)s!$ states.
Enumerate the histograms of $S$ independent end states on $1+s^2$
classes, at most $(S+1)^{1+s^2}$ possibilities, using the multinomial
law; this updates the vector or fails.  Compute the final end-state
distribution by the same dynamic program; apply the capped geometric sum
to the single-attempt distribution.  Enumerate legal matchings for $\pi$
only on the correction branch. Every partial law at a fixed clock time
has a common dyadic denominator of \emph{polynomial bit length}: pad
the entire attempt/retry horizon to its deterministic polynomial maximum,
with $b$ bits for a token and $b$ for acceptance at each transition.
Histogram multinomial computations likewise admit the common denominator
for $S$ padded samples.  This avoids taking coprime denominators over an
exponential number of random histories. The table work and scans, including
schoolbook arithmetic and a final exact cumulative rational selection, cost
$\exp(O(s^4\log^2(s+2)))$ bit operations in expectation. The asymptotic
coefficient in the log-width bound can be chosen independently of $C$:
$R=O(s^2\log s)$ and $\log S=O(\log s)+O(\log C)$, where the
latter additive constant contributes lower order as $s\to\infty$.
Any fixed power of the table width for straightforward scans costs the
same order in its logarithm. Take $C$ large enough that its prefactor
$2^{-u}$ pays for this whole branch.

On the fast branch, each of $O(1)$ changed edge and hole factors has
$O(R\log(S+1)+s\log(s+2))=\widetilde O(s^2)$ bits. Local probability
ratios therefore take $\widetilde O(s^4)$ schoolbook bit work.
Accessing a whole table of $s^2$ hole entries sequentially stays within
this bound. Do not construct the potentially long $b$-bit rounded
transition probability at every step. For a rational $a/d\in(0,1)$
reveal its binary digits by doubling and subtracting its remainder;
compare with fresh fair bits, stopping on the first discrepancy or rejecting
on a tie of $b$ digits. The expected number of digits used is at most two.
For $4s$ nearly uniform token bins, a dyadic prefix of length $k$ is
ambiguous only if it meets a boundary, probability $O(s2^{-k})$; the
expected prefix cost is polylogarithmic. These lazy operations have exactly
the stipulated finite dyadic law. A $2^{-u}$ correction coin likewise
short-circuits.  Thus $RST_0+T_1=\widetilde O(s^{15})$ transitions per
attempt, $O(s^2)$ expected attempts, and $\widetilde O(s^4)$ expected
bit work per transition give the asserted exponent. This is a fair-bit
adapter, not a relabeling of the arithmetic complexity in the source.
\end{proof}

\begin{proposition}[Full accounting]\label{prop:runtime}
The total expected fair-bit work, including all proposals and their
rejections and the one-time preparation, is $\widetilde O(n^{25})$.
\end{proposition}
\begin{proof}
We detail the interfaces that could otherwise conceal a further power
of $n$.  There are $O(n)$ feasible ordinary matching factors per history,
each of side at most $n$. For each designated fill of the balanced
lookahead block, use the free row and column sets of each color $k$,
with $v_k$ the size of its own pin set and $s_k=n-v_k\ge n/2$.
In the current free matching $F$ of size $s_c$, the number of assigned
edges free for the later color $k$ is $s_c-2v_k+h_k(F)$, where $h_k(F)$
counts assignments from the owner-$k$ domains to the owner-$k$ images.
The sets are disjoint and $\sum_{k>c}h_k(F)\le n$. Thus the weight is a
known rational times
\[
 \psi(F)=\prod_{k>c}(1-1/s_k)^{h_k(F)}\in[e^{-4},1]
\]
for large $n$. Rejection from an unweighted uniform matching has constant
expected cost. Its weighted normalizer is its known rational times
$\per(B)\E\psi(F)$. An independent exact $\E\psi(F)$-coin uses one
unweighted matching and one rational flip. The unbalanced fallback is
unweighted.  For each of the at most two cubic factors, write the capped
mark mass as $A\per(T)\theta_T$ with $A=n$,
$\theta_T=\E_M\min(1,2^{\#\text{cycles}(T-M)}/A)\ge1/A$.
Here too one unweighted matching gives each input coin, and a cubic
coloring draw costs at most $A$ such samples in expectation.

In Lemma~\ref{gate:binary-v2}, encode the anchor's next column by binary
splits of its current row. Conditional on all prior revelations, the
expected reciprocal of the selected branch probability is exactly the
number of positive children, at most two.  Use the deterministic upper
bound $N=O(n^2\log(n+2))$ on total sites, including the additional
scalar coins, and $m=N+2$ repetitions for the negative-binomial
estimators. Positive and reciprocal factors both use at most
$(m+1)/p$ expected probes at a site of probability $p$, also under the
specified size-biased simulation. The $O(n)$ ordinary lookahead scalar
coins have $1/p\le e^4$; the $O(1)$ cubic ones have $1/p\le n$.
Even those latter sites cost only $O(mn)$ total. The constant-mean
two-sided sigmoid race and the forward/reverse draws themselves are
therefore covered by $\widetilde O(n^4)$ unweighted matching calls per
soft gate, not $O(n^5)$ and not one weighted-graph call on a bigger board.

Here is generous accounting for everything else. In the fixed $3\times3$
cubic contingency table it suffices to index four entries and the layer
$t\le n$, since row and column sums are $t$; at most $O(n^5)$ states,
six predecessor transitions per state, and $O(n\log n)$-bit coefficients
before combining their squares and factorials. Even a full sequential
table scan per operation is below $O(n^{13})$ bit work. Each ordinary
marked-block composition table takes $O(bn^2)$ convolutions with
$O(bn\log n)$-bit entries after clearing by $n!$ at each step; the
physical block sizes have $\sum b\le n$. Fixed-length polygon and pair
cycle-index tables need only polynomially many factorials with
$O(n\log n)$ bits. Masks, feasibility via augmenting paths, row order
and shallow binary restrictions can all be scanned/constructed in
$\widetilde O(n^5)$ work outside individual probes. For rational
lookahead factors a crude bound on the aggregate known numerator and
denominator length is $O(n^4\log n)$: there are at most $n^2$ color
pairs, with exponents $O(n)$ on $O(\log n)$-bit factors. This includes
the factors on the realized matching. The logistic threshold carries
$O(n^2\log n)$ bits (and multiplication/division by the level index's
dyadic factor changes this by $O(\log n)$). Polynomially many applications
of exact rational arithmetic, with clearing of denominators as above,
fit comfortably in $O(n^{24})$ bit work. In particular, at most
$O(n^2)$ such lookahead factors can be evaluated by $O(\log n)$ repeated
squarings each. At a probe, evaluating a conditional scalar coin
directly costs at most $O(n^7)$ bit work, including a complete mask scan,
forming its exact rational, and sampling it; even this can be charged to
every probe without competing with $P(n)$.

The unbounded waiting counters are not presumed to fit in a fixed word.
Every positive binary branch has $p\ge1/n!$, and all added scalar
sites obey the bounds just given. For a negative-binomial waiting time
$V$ to $k\le m+1$ successes,
$\E[V\log(1+V)\mid p]\le(k/p)\log(1+(k+1)/p)$ by size-biased
Jensen. Uniformly its squared log length has expectation
$O((1+n\log n+\log m)^2)$ by Markov's inequality for $V$ at
exponential multiples of its mean. Consequently the $N$ unrestricted
counter factors' total bit length has second moment
$O(N^2(1+n\log n+\log m)^2)$. Multiply them without reducing and
combine with the deterministic prefactor at the end. These counter
updates, products and exact rational comparison are within the stated
$O(n^{24})$ allowance. Thus one forward or reverse soft unit has
expected cost $W(n)=\widetilde O(n^4P(n)+n^{24})$ uniformly over
inputs and dyadic levels.

Finally, the proved reverse mean bound and the incidence count lower
give $H/|\cL_n|=O(1)$ for a sufficiently large fixed multiple of the
rounded factorial scale. The explicit polynomial log-range bound $M$
in the multiscale flattening is sufficient: it contributes only
$O(\log^2(2M))$, not $M$ or $M^2$, to the expected work.
A polylogarithmically diverging safety padding of the scale also
leaves the soft exponent unchanged. The scale itself can be computed
once by factorials and outward log/exp intervals (additive error less
than one in the log); its $O(n^2\log n)$-bit length, and the
constant-degree polygon/contingency preprocessing, are below the same
overall bound. Substituting Lemma~\ref{lem:pmclock-v2} gives
$\widetilde O(n^{4+21})$. The finite small-order part can be handled
by drawing uniformly from all independent permutations and rejecting
non-Latin arrays; it does not alter the asymptotic statement. No
literal bound $O(n^{25})$, numerical crossover or worst-case
polynomial stopping time is implied by the displayed soft bound.
\end{proof}

\begin{proof}[Proof of Theorem~\ref{thm:main}]
For sufficiently large $n$, run the proposal of Sections~\ref{sec:coverage}--\ref{sec:lookahead-cubic} with the multiscale tests of Lemma~\ref{gate:multiscale-v2}, repeating independent trials until one returns a square. Proposition~\ref{prop:mean-cover}, Lemma~\ref{lem:midpoint} and Proposition~\ref{prop:count-lower} give both hypotheses \eqref{gate:hyp-v2} for a fixed upward-rounded multiple of $X_n$. The multiscale lemma gives a common strictly positive mass for each target in a trial, so repetition terminates almost surely and is exactly uniform. Lemma~\ref{gate:binary-v2} implements its tests with exact matching samples; Proposition~\ref{prop:runtime} bounds the total work. For the finitely many smaller orders, independent uniform permutation rows conditioned on column distinctness give an exact sampler by rejection. The cyclic Latin square shows positive acceptance at every such order.
\end{proof}

\section*{Acknowledgment}
The author provided the problem specification and research direction. Six OpenAI GPT-6.2 Astra research agents contributed to proof development, internal cross-checking, exposition and reference code.
\appendix
\section{A finite expansion and a local derivative for a regular permanent}
\label{an:appendix}
We give the analytic details with a \emph{fixed} positive degree exponent. Constants in this section can depend on that exponent and on a fixed Taylor order, but not on the board. This suffices even at the smaller cutoff $n^{1/200}$. Write $F=A/d$, $J=\mathbf 1\mathbf 1^{\mathsf T}/n$, $B=F-J$, $V_d=n!(d/n)^n$, $\ell=2+\lceil\log_2(2n)\rceil$, $\sigma=\|B\|_{\rm op}$ and $U=\|B\|_F^2=n/d-1$. All transposes in formal identities are transposes, not conjugate transposes. A mixed template is a finite connected bipartite multigraph, with each edge carrying kernel $F$ or $J$. Its value is the product of its kernels summed over every vertex label on its prescribed shore, without any inequalities on labels. Labels at different vertices are free even if their vertices belong to the same shore.

\begin{proposition}\label{an:expansion}
Fix $a>0$ and $m>0$. If $d\ge n^a$, $\sigma\le\ell^3/\sqrt d$, there is an explicitly generated fixed rational combination $P_A(z)$ of connected mixed templates times monomials in $z$, approximating the centered log rook polynomial as in (A6). There is another fixed combination $\Psi_A$ of the same kind (without $z$) for which
\[
 \log(\operatorname{per} A/V_d)
 =\Psi_A+\frac{U}{2n}
 +O_{a,m}\left(n^{-m}+\frac{\ell^{C_{a,m}}}{d^2}
                           +\frac{1}{d\sqrt n}\right).
 \tag{A1}
\]
In particular the last error is summable $o(1)$ over $d\ge n^a$ if $m>2$. The template lists, coefficients, and exponents depend only on $a,m$. For every allowed edge $e$, its probability $p_e$ in the \emph{uniform} perfect matching satisfies
\[
       p_e=d^{-1}+O_{a,m}(\ell^{C_{a,m}}d^{-5/2}+n^{-m}).
       \tag{A2}
\]
\end{proposition}

\paragraph{Rooks, a centered covariance and an absolutely convergent gas.}
Let $r_j(B)$ be the sum of the permanents of $j$-by-$j$ submatrices, and $R_B(z)=\sum_j r_j(B)z^j$. Classical rook expansion around $J$ gives
\[
 \operatorname{per}(F)/\operatorname{per}J
  =\sum_j\frac{n^j}{(n)_j}r_j(B)=\mathbb E R_B(1/X),
 \qquad X\sim\operatorname{Gamma}(n+1,\text{rate }n).
 \tag{A3}
\]
Indeed one chooses the rows and columns on which $B$ is used and the other entries all equal $1/n$; the negative moment of $X$ is $n^j/(n)_j$. Put
\[
 \Omega=\begin{pmatrix}0&B\\ B^{\mathsf T}&0\end{pmatrix},
 \qquad \Gamma=z\Omega(I+z\Omega)^{-1},\quad
 f(w)=(1+w)e^{-w+w^2/2},\quad g(w)=f(w)-1.
\]
The analytic Wick identity of Divoux--Wang \cite[Sec.~3]{DW}, in this notation, reads
\[
 R_B(z)=\det(I-z^2B^{\mathsf T}B)^{-1/2}\,
               Z,\qquad Z=\mathbb E\prod_{i=1}^{2n}f(\xi_i),
 \quad \mathbb E\xi_i\xi_j=\Gamma_{ij}.                    \tag{A4}
\]
The branch at zero is used. For completeness the formal covariance $z\Omega$ in $\mathbb E\prod(1+y_i)$ pairs a subset of lefts and rights, giving precisely the rooks. The contractions of $\sum_i y_i$ with every variable vanish since $\Omega\mathbf1=0$. On introducing $e^{-\sum y_i^2/2}$ against $\prod f(y_i)$, Gaussian square completion gives updated covariance $\Gamma$ and factor $\det(I+z\Omega)^{-1/2}$, as displayed. Wick convergence below justifies the formal operations for the small norms in use.

For $|z|\le 3/2$, $\|\Gamma\|_{\rm op}\le C\sigma$, $\max_{ij}|\Gamma_{ij}|\le C/d$, and both absolute row and column sums are at most $C_a$. To see the last assertion, each nonzero shore block of $\Omega^t$ is an alternating stochastic product of $F,F^{\mathsf T}$ minus $J$; its absolute row sums are at most two. Sum the first $\lceil 8/a\rceil+10$ powers by this fact and the remaining powers by the row-norm bound $\sqrt{n/d}\,\sigma^{t-1}$. For $t\ge2$ an individual entry costs $\sigma^{t-2}/d$ by the two endpoint Euclidean norms.

One may use extra independent circular Gaussian noise in (A4): holomorphic moments are determined by the complex symmetric pseudo-covariance $\Gamma$. Choose a realization with Hermitian covariance $2b I$ where $b>\|\Gamma\|_{\rm op}$ and $b=O(\sigma)$; the real covariance is
\[
 \tfrac12\begin{pmatrix}2b I+\Re\Gamma&\Im\Gamma\\
                    \Im\Gamma&2b I-\Re\Gamma\end{pmatrix}.
\]
Its eigenvalues lie in $[b/2,3b/2]$. These bounds remain true when site blocks of the real covariance are Schur multiplied by a positive semidefinite correlation matrix: the Schur map is positive and unital. The representation applies as well to symmetric, nonnormal complex pseudo-covariances. For $\sigma=0$ the relevant identities at the unperturbed matrix follow directly or by continuity.

For a set $S$ of distinct sites write $w(S)$ for the joint cumulant of $g(\xi_i),i\in S$. The moment--cumulant identity gives the \emph{hard-core} partition function $Z=\sum_{\mathcal C}\prod_{S\in\mathcal C}w(S)$ over unordered pairwise disjoint nonempty supports and the empty family. For $v=|S|\ge2$ its activity is exactly
\[
 w(S)=\sum_{\tau\ {\rm tree\ on}\ S}\prod_{ij\in\tau}\Gamma_{ij}
 \int_{[0,1]^{v-1}}\mathbb E\prod_{i\in S} g^{(\deg_\tau i)}(\eta_i)\,du .
 \tag{A5}
\]
In the forest interpolation formula \cite{AR}, let $S_{ij}(u)$ be the minimum parameter on the tree path between $i$ and $j$ (diagonal one); the covariances of $\eta$ are the original site blocks times $S_{ij}(u)$. This matrix is the integral of the component-equivalence matrices in thresholded forests and is therefore a correlation matrix. To prove (A5) for polynomials, the cumulant selects connected Wick pairings. Give the nonloop slot-pairs independent continuous uniform priorities. Mark their unique maximum spanning tree. Conditional on its priorities an omitted nonloop $ij$ must have priority below the minimum on its tree path, giving exactly $S_{ij}(u)$. Loops have factor one. Selecting the incident distinguished slots produces the derivatives at each site; tree neighbors distinguish the incident slots and there is no extra local factorial. Summing and integrating recovers each connected pairing once. For the entire $g$ and its derivatives use Taylor truncation: the absolute power-series majorants at each coordinate are bounded by $C_k e^{C|w|^2}$ at any fixed derivative order. Gaussian exponential integrability follows from the small operator norm for each finite system, uniformly in the weakening; this proves the limit including singular covariance cases.

Gaussian H\"older \cite[Thm.~1]{CDP} (the block covariance-domination form of Gaussian Brascamp--Lieb) applies with an absolute exponent at each site, by the real covariance bounds above. Cauchy's estimate at radius $\sqrt{k}$, together with the third-order zero at zero for $g$, gives
\[
 |g^{(k)}(w)|\le C^{k+1}\sqrt{k!}\,|w|^{(3-k)_+}e^{C|w|^2}.
\]
For $k\le2$ divide by the known zero before applying Cauchy; for $k\ge3$ use $k!k^{-k/2}e^{Ck}\le C'^k\sqrt{k!}$. Gaussian absolute moments at a site give $\|g^{(k)}(\eta_i)\|_{p_i}\le C^{k+1}\sqrt{k!}\,\sigma^{(3-k)_+/2}$. On a tree $\sum_i(3-\deg_\tau i)_+\ge v+2$. Sum labels outward fixing a prescribed root, dropping distinctness for this absolute bound. The row sums cost $C_a^{v-1}$. By Pr\"ufer coding, the sum of $\prod_i\sqrt{k_i!}$ over labeled trees divided by $(v-1)!$ is at most $[\sum_{j\ge0}\sqrt{(j+1)!}/j!]^v$. Thus the rooted total for all $|S|=v$, even with weight $e^{2v}$, is at most $(C_a\sqrt{\sigma})^v\sigma$.

There is a complementary bound for fixed $v$. For $E$ paired edges the sum of slot/valence multiplicities before labels is at most
\[
 (2E-1)!!\,[x^{2E}]\left(\sum_{h\ge3}|[w^h]g(w)|x^h\right)^v
 \le (Cv)^{C(E+1)}
\]
by evaluation at $x=\sqrt{E/v}$ and the Gaussian-type series majorant. A rooted spanning tree of the pairing graph costs $C_a^{v-1}$; each of the other $q=E-v+1$ edges costs $C/d$. Always $q\ge2$ and $v\le2(q-1)$, including a singleton whose nonzero Wick degree must be even and at least four.

\paragraph{Finite, enumerable connected truncation.}
Fix an integer $h$ sufficiently large in terms of $a,m$. Discard supports with more than $8h$ sites using (A5), and on smaller supports discard diagrams with $q\ge h$ using the fixed-site bound and the geometric tail in $E$. The rooted absolute errors, including $e^{2v}$, are at most $C_{a,h}\ell^{C_h}d^{-h}$. The full rooted sum is $o(1)$: for example separate the fixed small supports through $16$ with bound $O_a(d^{-2})$ and the larger tail using (A5). If $\sum_{S\ni i}|w(S)|e^{2|S|}<1/10$, removing a site gives $|Z_U/Z_{U\setminus\{i\}}-1|<1/5$ by induction on $|U|$: group the configurations by the unique support containing $i$ and use the earlier site ratios along its other sites. In particular $Z$ is nonzero and $|\partial\log Z/\partial w(S)|=|Z_{\bar S}/Z|\le e^{|S|}$. Interpolate linearly between exact and truncated activities to bound their log difference by $O_{a,h}(n\ell^{C_h}d^{-h})$.

Next weight each remaining diagram by $u^q$. On $|u|\le d/C_{a,h}$ the rooted weighted sum is still small, by the fixed-site bound and $q<h$. Thus the log of the gas there is $O(n)$ by the same induction. Its terms with total degree in $u$ at least $h$ cost $O_{a,h}(n/d^h)$ at $u=1$. For every retained monomial the support-intersection graph is connected: on setting all other activity indeterminates to zero disconnected groups factor, and their mixed log coefficient vanishes. Repeated occurrences of the same support are counted as overlapping. Enumerate the finite set of slot pairings and site-equality patterns using ordinary Taylor expansion of $\log(1+y)$ to degree $h-1$. Remove inequalities among the surviving labels by partition-lattice M\"obius inversion. The result consists of connected sums with at most $3h$ factors $\Gamma_{ij}$, allowing loops. This is a finite rational construction, not a summation over actual vertex labels by the sampler.

Replace each propagator by its first $M=2h+5$ powers of $z\Omega$. The error matrix has max entry $O(C^M\sigma^{M-1}/d)$ and operator norm $O(C^M\sigma^{M+1})$. For an error edge not a bridge choose a spanning tree avoiding it and use the max entry; for a bridge separate the two rooted sides, whose entry suprema are $O_{a,h}(1)$, and use the operator norm between the two vectors. In either case the free-root factor is at most $2n$ and the total error is $O_{a,h}(n\ell^{C_h}/d^h)$. This is done \emph{after} label inequalities are removed. Expand each power into alternating bipartite paths of $B$ and its transpose and each $B=F-J$; new internal vertices are unrestricted. The determinant half-log has terms $z^{2k}\operatorname{tr}[(B^{\mathsf T}B)^k]/(2k)$ for $1\le k<h$, with tail $O_h(U\sigma^{2h-2})$. These are connected cycles. We have constructed $P_A$ with
\[
 \sup_{|z|\le3/2}|\log R_B(z)-P_A(z)|\le C_{a,h} n\ell^{C_h}/d^h .
 \tag{A6}
\]
One may slightly enlarge the fixed disk first when invoking Cauchy at its boundary. Take $h$ such that this error is $O(n^{-m-4})$. Since $a>0$ this needs only a fixed $h$.

\paragraph{The real Gamma tilt.}
For complex zero-sum $B$, $|r_k(B)|\le\|B\|^k r_k(I-J)$: write $k!r_k(B)=\langle u,B^{\otimes k}u\rangle$ for the indicator $u$ of ordered injective $k$-tuples and project on both sides by $(I-J)^{\otimes k}$. Also $\sum_k n^k r_k(I-J)y^k/(n)_k=\operatorname{per}(J+y(I-J))/\operatorname{per}J\le e^{Cny^2}$ for $0\le y\le1/2$: expanding $(1-y)J+yI$ gives $(1-y)^n\sum_{k=0}^n(ny/(1-y))^k/k!$. Hence the signed Gamma mass outside a fixed small band $I$ about one is exponentially small: for $k<cn$ use the Gamma$(n+1-k,n)$ tail after reweighting by $X^{-k}$; for $k\ge cn$ extract $2^{-cn}$ in the tensor bound at $2\|B\|$. Here $\|B\|=o(1)$ and $c>0$ can be chosen smaller than the bandwidth. On real $I$, $R_B(1/x)>0$ by (A4) and zero-freeness connected to $z=0$; $\operatorname{per}F/\operatorname{per}J\ge1$ ensures the discarded absolute tail is negligible relatively. Replace $\log R_B(1/x)$ by $P_A(1/x)$ there.

Write $L(x)=P_A(1/x)$. The determinant's quadratic term is $U/(2x^2)$. Its higher terms have size $O(U\sigma^2)$ on a fixed larger band; the gas log is $O_a(n/d^2+n(C_a\sqrt{\sigma})^{17}\sigma)$. Consequently $L(x)=U/(2x^2)+E(x)$ with $\sup_I|E^{(j)}|\le C_{a,h}n\ell^{C_h}/d^2$, $j\le3$, by Cauchy on a smaller fixed band. For formal scalar $s$ let
\[
 t=\frac{sP_A'(1/(1-t))}{n(1-t)},\qquad
 \Phi(s)=n\{\log(1-t)+t\}+sP_A(1/(1-t)).
\]
On $|s|\le d/C_{a,h}$ this is a contraction on a fixed small disk in $t$, since $|P_A'|+|P_A''|=O_a(n/d)$ there. $\Phi$ is analytic and bounded by $O(n)$, so Taylor through order $h-1$ at $s=1$ costs $O_{a,h}(n/d^h)$. The coefficients are obtained by Lagrange inversion or truncated rational power-series arithmetic: $\Phi'(s)=P_A(1/(1-t(s)))$ and, for $j\ge2$,
\[
 [s^j]\Phi=\frac{[y^{j-2}](P_A'(1/(1-y)))^j(1-y)^{-j-1}}{j(j-1)n^{j-1}} .
\]
Join $j$ connected graphs by $j-1$ constant $J$ edges, each joining named opposite-shore vertices of consecutive graphs, to implement $n^{1-j}$ exactly without identifying vertices. This constructs $\Psi_A$.

At $s=1$ put $x_* =1-t=1-U/n+O(\ell^{C_h}/d^2)$. Changing rate to $n/x_*$ gives correction, beyond $e^{\Phi(1)}$,
\[
 x_*\,\mathbb E[ e^{G(Y)}\mathbf1_{\{x_*Y\in I\}}],\quad
 Y\sim\operatorname{Gamma}(n+1,n),\quad
 G(Y)=L(x_*Y)-L(x_*)-x_*L'(x_*)(Y-1).
\]
Here $L'(x_*)=n(1-1/x_*)$. On the retained band $|G|\le C(n/d)(Y-1)^2$. Its mean there equals $\tfrac12 x_*^2 L''(x_*)(1/n+2/n^2)+O(1/(d\sqrt n))$ since the centered third absolute moment is $O(n^{-3/2})$. The log moment differs from its mean by $O(1/d^2)$, using $\mathbb E[(Y-1)^4 e^{C(n/d)(Y-1)^2}\mathbf1_{\{x_*Y\in I\}}]\le C/n^2$ from Gaussian tail bounds for the Gamma density \emph{on this band}. Tails of the ordinary moments are exponentially small. Thus $\log x_*=-U/n+O(\ell^{C_h}/d^2)$ and the curvature contributes $3U/(2n)+O(\ell^{C_h}/d^2+1/(d\sqrt n))$. This gives (A1), including the sign of its common correction. One cannot charge an unspecified graph-dependent $O(1/d)$ in a comparison of two path laws requiring $o(1)$ total error.

\paragraph{A balanced tangent and a rooted derivative.}
For allowed edge $e=(i,j)$ put $r=e_i/d$, $c=e_j/d$ as column vectors, and solve on the mean-zero subspace
\[
 u=(I-FF^{\mathsf T})^{-1}(r-Fc),\quad v=c-F^{\mathsf T}u,\qquad
 T_{xy}=F_{xy}(\mathbf1_{\{(x,y)=e\}}-u_x-v_y)
\]
for all row and column indices $x,y$. Then $T$ has zero row/column sums, is supported on $A$, and $\|T\|_{\rm entry,1}:=\sum_{ij}|T_{ij}|\le C_a/d$. Indeed bound the first $\lceil4/a\rceil+10$ Neumann powers in $\ell^1$ by stochasticity and the remaining ones by $\sqrt n\sigma^{2k}$ on the centered right-hand side. Summing the potentials along an arbitrary perfect matching and using $\sum u+\sum v=1/d$ gives
\[
 \left.\partial_\theta\log\operatorname{per}(F+\theta T)\right|_0=p_e-1/d,\qquad
 \langle B,T\rangle=0 .
\]
The latter equality holds since $B$ is constant on the support and $\sum T_{ij}=0$.

Apply (A4) to this affine perturbation for $|\theta|\le n^{-6}$ with complex coefficients. Both matrices remain centered after subtracting $J$; absolute row/column sums of $F_\theta$ are at most $1+O(n^{-6})$. In the earlier bounds use $\rho=2\ell^3/\sqrt d$ in place of $\sigma$ and the same initial finite number of alternating powers. The isotropic-noise representation above proves holomorphy of all site moments by Wick and Taylor convergence for these complex perturbations; it is not necessary to differentiate a matrix absolute value or the law of a nonholomorphic covariance realization.

The derivative of the determinant half-log is
$\sum_{k\ge1}z^{2k}\langle B(B^{\mathsf T}B)^{k-1},T\rangle$.
Its first term vanishes. The two endpoint Euclidean norms for the others bound each entry by $\sigma^{2k-3}/d$, hence total $O(\sigma\|T\|_{\rm entry,1}/d)$. For the gas use
$\partial\log Z=\sum_S (\partial w(S)) Z_{\bar S}/Z$. The rooted absolute activity on supports larger than a fixed $v_0=v_0(a,m)$ costs $O(n^{-m-10})$ even after summing over all roots, by (A5) uniformly on $|\theta|=n^{-6}$. Cauchy under the contour integral gives derivative tail $O(n^{-m-4})$. For smaller supports differentiate the absolutely convergent Wick pairings. The distinguished factor is
\[
 \Gamma'=(I+z\Omega)^{-1}z\Omega'(I+z\Omega)^{-1},
 \quad \|\Gamma'\|_{\rm entry,1}\le C_a\|T\|_{\rm entry,1},
 \quad \sum_i|\Gamma'_{ii}|\le (C_a/d)\|T\|_{\rm entry,1}.
\]
For the diagonal bound expand $(I+z\Omega)^{-1}=I-\Gamma$ on both sides. $\Omega'$ has zero diagonal; any nonzero diagonal convolution costs one maximum entry $C/d$ and, if necessary, one absolute row sum $C_a$. For a distinguished nonloop use a spanning tree containing it: sum both endpoints using the total entry norm, and all $q=E-v+1\ge2$ remaining chords cost $C/d$. For a distinguished loop the root sum uses the diagonal bound and the other $q-1$ chords provide the remaining powers. The fixed-$v$ slot bound above thus sums, including the $E$ choices for the differentiated pair, to $O_{a,m}(\|T\|_{\rm entry,1}/d^2)$. With the hard-core derivative bound this controls the log gas uniformly on the central band.

Do \emph{not} differentiate an uncontrolled saddle remainder. Integrate the log-rook derivative on the positive real central Gamma band, dividing by $\operatorname{per}F/\operatorname{per}J\ge1$. The exterior integral is holomorphic in $\theta$ and bounded in absolute value by $e^{-cn}$ on its $n^{-6}$-circle by the tensor projection argument, so contributes at most $n^6e^{-cn}$ to the derivative. Combining the bounds just obtained proves (A2).

\paragraph{From one edge to a common template reference.}
We record details used at two different places in the counting and coverage proofs. In a healthy $s$-regular board at $s\ge n^a$, for any bounded list of already pinned disjoint edges, the conditional probability for one more allowed disjoint edge is $(1+O_a(1/s))/s$ (only its upper bound is needed after the second pin). In the matching-contracted graph a loopless transition matrix is $(A M^{-1}-I)/(s-1)$, centered norm $\sigma'=O(\ell^3/\sqrt s)$, maximum entry $(s-1)^{-1}$. Powers of a sufficiently large \emph{constant} length have all entries $(1+O(n^{-10}))/n$. To insert the missing edge, count simple paths back from its column-owner to its row of length $h-1$ and switch the closed alternating cycle; to remove it count length-$h$ simple cycles rooted at its row. For an internal repetition with shorter loop of length $j\ge2$ its probability relative to the conditioned endpoint is $O_a(1/s+1/n)$: bound the loop diagonal by $1/n+\sigma'^{j-2}/(s-1)$ and use mixing on the longer exterior chain. For a long loop use its $O(1/n)$ diagonal and the remaining short path bound $O(1/s+1/n)$, excluding the forced endpoint coincidence. Hits to a pinned row are bounded by mixing one side and $O(1/s)$ for the other; neither endpoint is a pinned row. There are $O_a(1)$ positions. This gives the asserted switch ratio. Conditional on distinct specified cells still present, (A2) therefore applies to each unconditioned singleton, joint disjoint pairs have probability $s^{-2}(1+O_a(s^{-1}+\ell^C/s^{3/2}))$, and triples have probability $O_a(s^{-3})$.

If $H$ is any fixed injective simple pattern with $e$ edges, set $\alpha=\sum_v\binom{\deg_H v}{2}$. Three-term Bonferroni for removal of the next matching gives survival, relative to $(1-1/s)^e$, of
$1-\alpha/s^2+O_{H,a}(\ell^C/s^{5/2}+n^{-m})$. Iterating from the full board gives the normalized mean
\[
 \frac{\mathbb E\operatorname{inj}(H,A_d)}{d^e}
 =\frac{(n)_{v_L}(n)_{v_R}}{n^e}
 \left[1-\alpha(1/d-1/n)+O_{a,H}(\ell^C/d^{3/2})\right]
 +\text{negligible error}.
\]
For fixed $H$ choose the spectral-failure exponent greater than $2e+2m+20$. Spectral goodness for the decreasing chain follows by stopping the bounded-increment martingale for $A_s/s$: conditional expectation of its drift has operator norm $O_a(1/s^2)$ while the coarser switch estimate holds on a healthy board (when only $\sigma\le1/2$ is known, use $O(\ell)$ switch length and $O(\ell^2/s^2)$ drift). The increment squares sum to $O(1/d)$, so matrix bounded differences plus union over $d\ge n^a$ ensure the claimed $\sigma$ bound and prevent exit. For a random $d$-subset of any fixed target's row permutations, matrix Bernstein on independent inclusion at rate $d/n$ followed by conditioning on subset size $d$ (cost $\le n+1$) gives the same bound. Only fixed patterns occur here.

Removing all $J$-edges of a connected mixed template gives $c$ components with $c\le a_0+1$ if $a_0$ such edges were removed. Retain the original $d^{-j}$ even if parallel $A$ edges are suppressed. Split assignments by exact partitions on each shore; whenever an injective forest occurs, use its unrestricted hom sum $n^{c'}d^{e'}$ minus nonidentity quotients. Quotient chains shorten strictly. Thus all terms are exact forests or $n^{-a_0}d^{-j}\operatorname{inj}(H,A)$ for cyclic simple $H$. On the latter $j\ge e$, $n^{-a_0}(n)_{v_L}(n)_{v_R}/n^e\le1$ since $e-v+c\ge1$. Inserting the preceding normalized mean gives the polynomial-time computable signed-template scalar used for $\mu_{\rm hi}$; determinant, true Gamma correction and $V_d$ factors are not discarded.

For the corresponding row subset of one \emph{fixed} square, write $P_s=(d)_s/(n)_s$, $N_H=(n)_{v_L}(n)_{v_R}$, $t=e-b$ for an equality partition $\rho$ of the row names on its edges using $b$ names, $M_\rho$ its number of labeled embeddings, and $m_2(\rho)=\sum_{C\in\rho}\binom{|C|}{2}$. Exactly
\[
 \mathbb E\operatorname{inj}(H,A_d)=
 N_H\left[P_e+\frac{\binom e2-\alpha}{n-1}(P_{e-1}-P_e)\right]
 +\sum_{t\ge2}\bigl(P_b-P_e-m_2(\rho)(P_{e-1}-P_e)\bigr)M_\rho.
\]
A pair of disjoint edges in a uniformly injected placement has equal names with probability exactly $1/(n-1)$ (fix both left images and one right image); adjacent pairs never do. This proves the identity including both signs. Dividing the bracket by $(d/n)^e$ gives the same $1-\alpha(1/d-1/n)+O_H(d^{-2})$. For each residual term, of either sign, the absolute multiplier is at most $(2+e^2)P_b$. Normalize and sum over $d>D$: if $X_\rho$ is the number on $b$ uniformly assigned distinct colors, the cost is at most $O_H(D^{-1})n^{-a_0}\mathbb E X_\rho$, since $M_\rho=(n)_b\mathbb E X_\rho$ and $j-b\ge2$. A connected colored component extends from a chosen root uniquely if at all. Choose a simple cycle in a cyclic component: all other components cost at most $n^{c-1}$ and its own component is determined by the rooted cycle. Thus this is at most $O_H(D^{-1})\mathbb E X_C$ for the selected colored cycle. A singleton color on $C$ bounds its expected count by $n/(n-b+1)$ by conditioning on the other colors and solving its closing edge. Otherwise every color on $C$ repeats; the finite polygon marking construction in the main text pays it. There are only a fixed number of patterns, so these costs can be aggregated even on the same word. This comparison concerns \emph{signed means} of the actual logarithmic score, not independent Poisson assumptions about overlapping motifs.

\section{The number of ordered factorizations: incidence without replacement}
\label{sec:incidence}
For this section write $U_d$ for the uniform law on the $N_d$ simple
$d$-regular boards on the two prescribed labelled shores, $W(A)=\operatorname{per}A$,
$m_d=\mathbb E_{U_d}W$, and $\widehat U_d(A)=W(A)U_d(A)/m_d$ for the
perfect-matching-biased law. Write $t_d=\|U_d-\widehat U_d\|_{\rm TV}$
and $V_d=n!(d/n)^n$. These probability spaces are used only to count:
the sampling algorithm never needs a uniform-board or a permanent-counting oracle.
We shall prove a constant-factor lower bound, rather than a power-of-$n$
lower bound which would change the running-time exponent.

\begin{proposition}[Constant-loss count lower]\label{prop:count-lower}
For an absolute $c_*>0$ and all sufficiently large $n$,
\[
 |\mathcal L_n|\ge c_*\prod_{d=1}^n m_d .
\]
\end{proposition}

\paragraph{Killing a marked incidence.}
Counting incidences $(A,M)$ with $A$ a $d$-board and $M$ one of its
matchings gives $N_d m_d=N_{d-1}m_{n-d+1}$, by deleting $M$.
At $A$ choose a uniform matching $M$, put $T=A-M$, and retain the step
with probability
\[
 \min\{1,W(A)/m_d\}\min\{1,m_{n-d+1}/W(K_{n,n}-T)\}.
\]
Starting from any subprobability dominated by $U_d$, the first factor
discards mass at most $t_d$, and the surviving marked incidence is
dominated pointwise by the uniform incidence. Its $T$ marginal is
therefore bounded by $U_{d-1}(T)W(K_{n,n}-T)/m_{n-d+1}$.
The second factor discards at most $t_{n-d+1}$ and restores domination
by $U_{d-1}$. Each particular matching transition now has probability
at most $1/m_d$. We need the following three estimates:
for fixed large $c$, with $D_0=\lfloor n^{1/10}\rfloor$ and
$R=\lceil\sqrt n\,\ell^{B_0}\rceil$ for a sufficiently large fixed $B_0$,
\[
 \sum_{c<d\le D_0}t_d\le C/c+o(1),\qquad
 \sum_{D_0<d<R}t_d=o(1),\qquad \sum_{R\le d\le n}t_d=o(1).
\]
The cutoff $D_0$ is internal to this incidence proof.
\paragraph{Sparse strip.}
For $c\le d\le D_0$, enumeration of simple bipartite boards with
prescribed sparse margins and forbidden cells gives
\begin{align}
 m_d&=n!d^{2n}\frac{((d-1)n)!}{(nd)!}
                      \exp[-1/2+O(d^3/n)],\notag\\[-1mm]
 t_d&\le C\{d^{-2}+d^{3/2}/\sqrt n+(nd)^{-1/4}\}.\label{eq:inc-sparse}
\end{align}
For details, in the formulas of Greenhill et al.\ \cite[Theorems~2.1--2.2]{GHIM}, let
$k_i,l_j\in\{d,d-1,d-2\}$, $s=\sum k_i=\sum l_j=\Theta(nd)$,
$S_2=\sum(k_i)_2$, $T_2=\sum(l_j)_2$, and
$F_X=\sum_{ij\in X} k_i l_j$ for a forbidden board $X$ of maximum
degree two. Their enumeration in this range reduces uniformly to
\[
 \frac{s!}{\prod k_i!\prod l_j!}
 \exp[-S_2T_2/(2s^2)-F_X/s+O(d^3/n)].
\]
Indeed $d^2=o(s^{2/3})$, their forbidden parameter is $O(d^2)=o(s)$,
$F_X=O(nd^2)=o(s^{5/3})$, and $O(d^2)F_X=o(s^2)$.
For one compulsory matching apply the formula to its complement after
decrementing its endpoint degrees by one. For two specified matchings
with $j$ common edges set $s=n(d-2)+j$ and $u=(d-1)j/s$.
Directly $S_2/s=T_2/s=d-3+u$ and
$F_X/s=2(d-2)-(d-3)u$, giving a net exponential relative to the
unrestricted count of $\exp(-u^2/2+O(d^3/n))$. There are
$(n!)^2 h_{n-j}/j!$ pairs where $h_m$ is the fraction of derangements
of $m$ objects. Set $x=(d-1)^{-2}$, $s_0=n(d-2)$ and let $K$ have
mass $\binom{s_0+j}{j}(1-x)^{s_0+1}x^j$. Extracting the factorials
and using Stirling yields
\[
 \mathbb E_{U_d}(\operatorname{per}A)^2/m_d^2
 =(1-x)^{-1/2}(1+O(d^3/n))
 \mathbb E[\mathbf1_{K\le n}e h_{n-K}e^{-u(K)^2/2}].
\]
Here $\mathbb EK=n/d+O(1)$, $\operatorname{Var}K=O(n/d)$, $K>n/2$
has exponentially small probability, $u(n/d)=1/(d-1)$ and
$|u'|=O(1/n)$. The expectation is $e^{-x/2}+O(1/\sqrt{nd})$.
Since $(1-x)^{-1/2}e^{-x/2}=1+O(x^2)$, total variation is bounded
by half the coefficient of variation, as required. Stirling in the
first line gives
$\log(m_d/(L_d\sqrt{n/d}))=O(d^{-1}+d^3/n)$ also for $2\le d<c$:
only the one-mark enumeration is used there. At $d=1$, $m_1=L_1=1$.

\paragraph{Moderate strip.}
Take $D_0<d<R=\lceil\sqrt n\,\ell^{B_0}\rceil$, $B_0$ fixed sufficiently
large, and $p=d/n$, $\tau=1/(nd^{0.9})+d/n^2+p^3$.
The bipartite case of Liebenau--Wormald \cite[Theorem~1.4]{LW} implies that for any fixed embedded simple pattern $F$
with $h$ edges, degrees $h_v$, and $w(F)=\sum_v\binom{h_v}{2}$,
\[
 \Pr_{U_d}(F\subset A)
 =\frac{\prod_v(d)_{h_v}}{(nd)^h}\exp\{w(F)/n+O_F(\tau)\}.
 \tag{I1}\label{eq:inc-pins}
\]
This is not a direct invocation with forbidden cells. After any bounded
number of margin decrements the unrestricted edge marginal is
$q_{ij}=(s_i t_j/m)(1+O_F(1/(nd^{0.9})))$: in that theorem take
$\kappa=.01$, $m=nd-O(1)$, endpoint deviations $O(1)$ and margin
variance $O(1/n)$, at density $o(1)$. Decrement the edges successively
using the exact ratio of margin counts $q_e^{\rm before}/(1-q_e^{\rm after})$,
then multiply by the probability of avoiding $F$ in the final reduced
law. In that law any two fixed pins have probability
$p^2(1+O_F(1/d+p))$ and three cost $O_F(p^3)$ by the same ratios
and union bounds. Thus log avoidance is
$-\sum_{e\in F}q_e^{\rm final}-hp^2/2+O_F(p^3+p^2/d)$.
The inverse denominators cancel the quadratic term.
In $\sum(q_e^{\rm after}-q_e^{\rm final})$ each later shared endpoint
contributes $1/n$, giving $w(F)/n+O_F(\tau)$. In particular
disjoint patterns factor to relative $1+O_F(\tau)$.

For connected unicyclic $F$ let $Y_F=\operatorname{inj}(F,A)/d^h$.
Its mean is $1-w(F)/d+O_F(d^{-2}+n^{-1}+\tau)$ and its variance is
$O_F(\tau+(nd)^{-1}+d^{-4})$. Indeed disjoint placements use \eqref{eq:inc-pins};
shared vertices but no edges change the normalized weight by
$O_F(1/d)$ on an $O_F(1/n)$ fraction (merge the vertex falling
factorials). For $j\ge1$ shared edges on $u$ shared vertices the
joint bound is $O_F(n^{j-u}/d^j)$; either $u\ge j+1$, or the
common graph contains the cycle, hence $j\ge4$ and $u\ge j$.
The absolute mean deviation from $m_F=1-w(F)/d$ is summable $o(1)$
over this strip, since in particular $\sum\sqrt\tau=o(1)$.

For $\widehat U_d$, delete the distinguished matching to obtain a $(d-1)$-board
tilted by its number of \emph{complement} matchings, then add its
conditional uniform complement matching. The tilt density is at most
$2$ by Br\'egman/van der Waerden since $p\log n=o(1)$.
For fixed $j$ disjoint available pins the complement matching probability
is $n^{-j}(1+O_j(p+1/n))$: switching two matched rows gives $n-O(d)$
available switches out of a pin; the unique reverse insertion fails
only if a forbidden cross-edge is present, with probability $O(d/n)$
by the two-pin upper $O(n^{-2})$ (the same upper applies after deleting
a matched row and column). The part of $Y_F$ with no new edge is thus
within summable absolute-mean error of $1-[w(F)+h]/d$.
For one specified edge marked, let $Y_e$ be its contribution. Under
the untilted remainder the first two moments of $dY_e$ give
$\mathbb E(dY_e-1)^2=O_F(1/d+p)$. To see this, disjoint embeddings in
the second moment contribute at most $1+O_F(1/d+p)$ by \eqref{eq:inc-pins} and the
two-mark bound. An overlap on $u$ vertices, $j$ old edges and $k=0,1$
common marks costs $O_F(n^{j+k-u}/d^j)$: if cyclic, $j\ge3$, otherwise
$u\ge j+k+1$. The first moment is $1+O_F(1/d+p)$ by the one-mark
estimate and a conditional absence probability $1-O_F(p)$ on the marked
cell. The nonnegative squared-error bound survives the bounded tilt.
Hence each singleton mark adds $1/d$ with summable absolute-mean error
$O_F(d^{-1}\sqrt{1/d+p})$. Terms with two or more marks cost
$O_F(d^{-2})$. This proves the same $m_F$ centering under both laws.

Apply the finite connected-template expansion of the log rook polynomial
before its Gamma averaging, at a fixed order $s>30$ (enlarge the order
if necessary so $\sum n\ell^{O(s)}/d^s=o(1)$). In each template delete
the $j$ constant $J$ edges; if $h'$ components remain, $j\ge h'-1$.
Expand homomorphisms into collision partitions, retaining tree
homomorphism terms exactly ($nd^e$), and cyclic simple injections.
A single unicyclic injection is replaced by $m_F$ after normalizing
by its number of distinct edges; retain the original further powers
of $1/d$ for repeated edges. Rank at least two costs $O_F(1/n)$ in
absolute mean by \eqref{eq:inc-pins} or the planted-mark decomposition. With $w\ge2$
cyclic components the remaining tree factors supply at most $n^{1-w}$;
two normalized spanning unicyclic counts have bounded joint mean
by the overlap argument above, and each further one is at most $n/d$.
Thus these terms cost $O_F(1/n)$. Only finitely many templates are
involved. Let $P_0$ be the resulting deterministic polynomial and set
\[
 c_d=\log\mathbb E[\mathbf1_{X\in I}\exp P_0(1/X)],\qquad
 X\sim\operatorname{Gamma}(n+1,n),\quad I=[1-1/(2\ell),1+1/(2\ell)]
\]
(shape, rate). On healthy boards the finite polynomial $P_A$ approximates
the actual log rook polynomial on the inverse window; the discarded
signed Gamma integral is $e^{-\Omega(n/\ell^2)}$ in absolute value,
negligible in log since $\operatorname{per}A/V_d\ge1$.
Theorem~B of Tikhomirov and Youssef \cite{TY} for the directed
loops-allowed model, gives the required centered singular bound with
arbitrarily high polynomial confidence under $U_d$; the planted
remainder density bound transfers it to $\widehat U_d$. Bad events have
polynomial log cost. Sandwiching the two positive integrals by
$\exp(\pm\sup_I|P_A(1/x)-P_0(1/x)|)$ gives
\[
 \sum_{D_0<d<R}(\mathbb E_{U_d}+\mathbb E_{\widehat U_d})
     |\log(\operatorname{per}A/V_d)-c_d|=o(1).\tag{I2}\label{eq:inc-common}
\]
This log argument does not require positive exponential moment control
of the error on the exceptional boards.

\paragraph{Dense strip.}
For $R\le d\le n$ let $z=d(d-1)/(n-1)$,
$c_{ij}=|N(i)\cap N(j)|$,
$W=\sum_{i<j}(c_{ij}-z)^2$, and put
\[
 w_*=(n/2)(d-z)(d-1)(1-(d-1)/(n-2)).
\]
Under $U_d$ or $U_d$ conditional on a fixed planted matching, resplit
the exclusive columns of a uniformly chosen row pair into equal halves
(subject to the at most two exclusive pins). This gives an exchangeable
pair $W,W'$. Let $2t$ be the exclusive-set size, $b_k$ its overlap
with row $k$, $D_k=c_{ik}-c_{jk}$. Then
$W'-W=\sum_{k\ne i,j}((D_k')^2-D_k^2)/2$; the subtracted half
averaged over pairs is $2W/n$. Unpinned,
$\mathbb E(D'_k)^2=b_k(2t-b_k)/(2t-1)$. Its sum differs from
$2t(d-1)(1-(d-1)/(n-2))$ by $O(d)$ on the good spectral event:
the squared deviations of the $b_k$ from their mean total $O(dt)$
by applying the centered board to the exclusive-set indicator.
Forcing up to two columns changes the sum by $O(d)$; for $t>2$ the
hypergeometric conditional variance and squared-mean changes per row
are $O(r_k+b_k/t)$, where $r_k$ counts pins incident to row $k$.
Smaller $t$ is immediate. Thus the drift is $-2(W-w_*)/n+O(d)$.
With both endpoints spectrally healthy, $|W'-W|=O(d^2)$ by applying
the centered norm bound to the difference of the pair's row indicators.
Bad endpoints have arbitrarily small polynomial mass under either law:
for dense $d$ use complementing, and transfer the plant by bounded
polynomial Br\'egman density (for $d<n/2$ via its remainder).
The identity
$-\mathbb E[(W-w_*)(W'-W)]=\mathbb E(W'-W)^2/2$
now yields $\mathbb E|W-w_*|=O(nd+\sqrt{n}\,d^2)$ under both laws.

The quartic signed log rook expansion has only one graph-dependent
term on the Gamma window, namely $x^{-4}W/(2d^4)$. Indeed, with
$B_A=A/d-J$, the trace quartic equals
$[nd^2+2W+n(n-1)z^2]/d^4-1$. Its remainder is
$O(n\ell^2/d^3+n^{-20})$ on healthy uniform or planted boards.
Replacing $W$ by $w_*$ in the truncated positive integral gives \eqref{eq:inc-common}
also over $R\le d\le n$ with summed bound
$O(\ell^{2-2B_0}+\ell^{-B_0})+o(1)$.


\paragraph{From the three strips to total variation.}
Write $h=\log(W(A)/V_d)$ and $M_d=\log(m_d/V_d)$.
Relative entropy in both directions gives
$\mathbb E_{U_d}h\le M_d\le\mathbb E_{\widehat U_d}h$.
Since $|e^x-1|\le(1+e^x)|x|$,
$2t_d\le(\mathbb E_{U_d}+\mathbb E_{\widehat U_d})|h-M_d|$.
The common centers in \eqref{eq:inc-common} thus prove the two upper-strip
claims; \eqref{eq:inc-sparse} proves the first since
$\sum_{d\le D_0}d^{3/2}/\sqrt n=O(n^{-1/4})$ and
$\sum_{d\le D_0}(nd)^{-1/4}=O(n^{-7/40})$.
In particular, for $c$ a sufficiently large \emph{fixed} integer, the
middle experiment starting at $U_{n-c}$ and stopping at degree $c$ retains
at least $1-2\sum_{d=c+1}^{n-c}t_d\ge 3/4$ for large $n$. Both
endpoint marginals of the surviving measure are dominated by their
respective uniform laws. No independence of the endpoints is asserted.

\paragraph{Why neither endpoint loses a power of $n$.}
Let $F(G)$ be the number of ordered $c$-factorizations of a $c$-board,
let $R_c$ be the number of $c$-row Latin rectangles, and set
$\mu_c=R_c/N_c=\mathbb E_{U_c}F$. Drawing $c$ independent uniform
perfect matchings conditional on disjointness and forgetting their order
gives exactly the density $F/\mu_c$ with respect to $U_c$. For fixed
$c\ge3$, this law and $U_c$ are mutually contiguous:
the graph-restricted sum of $c$ one-factors on prescribed labeled shores has this contiguity by Greenhill, Kim and Wormald \cite[proof of Corollary~1.1, equation~(1)]{GKW}. It is the perfect-matching contiguity in their proof, not their result on Hamilton cycles, that is used here.
Consequently for any $\varepsilon>0$ some $\delta>0$ satisfies
$U_c(F<\delta\mu_c)<\varepsilon$ eventually. Otherwise choose a
subsequence with $U_c(F<\mu_c/j)\ge\varepsilon$. The tilted masses
on it are at most $1/j$; extend the events by the empty set off the
subsequence to contradict contiguity.

Restrict the surviving middle measure to paths whose initial
\emph{complement} and final board both have $F\ge\delta\mu_c$.
With $\varepsilon=1/8$ the two removals cost at most $2\varepsilon$
even jointly, by endpoint domination. Every particular surviving
middle path (including its initial state) has probability at most
$N_c^{-1}\prod_{d=c+1}^{n-c}m_d^{-1}$, whereas every retained path
extends in at least $(\delta\mu_c)^2$ ways by ordering top and bottom
factors. A labelled complete square determines its middle path and
both endpoint orders uniquely. We obtain
\[
 |\mathcal L_n|\ge \tfrac12\delta^2\frac{R_c^2}{N_c}
                   \prod_{d=c+1}^{n-c}m_d .
\]
Finally $R_c=(1+o(1))\prod_{d=n-c+1}^n m_d$: at each of the first
$c$ steps the ratio of the Br\'egman and van der Waerden bounds on the
remaining $(n-O(c))$-board is $\exp(O_c(\log n/n))$ by Stirling,
so every prefix has the same extension factor to this accuracy.
The incidence identities telescope exactly to
$\prod_{d=n-c+1}^n m_d=N_c\prod_{d=1}^c m_d$.
Substitution proves Proposition~\ref{prop:count-lower}. The constants
in contiguity need not define a computable membership test: the two
endpoint cuts appear only in this counting argument.

\begin{thebibliography}{99}
\bibitem{JM}
M.~T. Jacobson and P. Matthews,
Generating uniformly distributed random Latin squares,
\emph{J. Combin. Des.} \textbf{4} (1996), 405--437.
\bibitem{GLMP}
A. G\"obel, J. Liu, P. Manurangsi and M. Pappik,
Perfect sampling from rapidly mixing Markov chains, Theorem~2,
\href{https://arxiv.org/abs/2410.00882v2}{arXiv:2410.00882v2} (2024).
\bibitem{Gurvits}
L. Gurvits, Van der Waerden/Schrijver--Valiant like conjectures and stable
(homogeneous) polynomials: one theorem for all,
\emph{Electron. J. Combin.} \textbf{15} (2008), R66.
\bibitem{LS}
M. Laurent and A. Schrijver, On Leonid Gurvits's proof for permanents,
\emph{Amer. Math. Monthly} \textbf{117} (2010), 903--911, esp.\ (6).
\href{https://doi.org/10.4169/000298910X523380}{doi:10.4169/000298910X523380}.
\bibitem{DW}
A. Divoux and S. Wang, The maximum number of Hamiltonian cycles in
tournaments, Section~3, \href{https://arxiv.org/abs/2610.23755v1}{arXiv:2610.23755v1} (2026).
\bibitem{GHIM}
C. Greenhill, M. Hasheminezhad, I. Iliffe and B.~D. McKay,
Asymptotic enumeration of constrained bipartite, directed and oriented graphs
by degree sequence, Theorems~2.1--2.2,
\href{https://arxiv.org/abs/2601.04822v1}{arXiv:2601.04822v1} (2026).
\bibitem{LW}
A. Liebenau and N. Wormald, Asymptotic enumeration of digraphs and bipartite graphs
by degree sequence, \emph{Random Structures Algorithms} (2023), Theorem~1.4.
\href{https://doi.org/10.1002/rsa.21105}{doi:10.1002/rsa.21105}.
\bibitem{Huber}
M. Huber, Nearly optimal Bernoulli factories for linear functions,
\emph{Combin. Probab. Comput.} \textbf{25} (2016), 577--591.
\href{https://arxiv.org/abs/1308.1562}{arXiv:1308.1562}.
\bibitem{Tropp}
J.~A. Tropp, User-friendly tail bounds for sums of random matrices,
\emph{Found. Comput. Math.} \textbf{12} (2012), 389--434;
Corollary~7.5 (matrix Azuma).
\href{https://doi.org/10.1007/s10208-011-9099-z}{doi:10.1007/s10208-011-9099-z}.
\bibitem{AR}
A. Abdesselam and V. Rivasseau, Trees, forests and jungles: a botanical garden for cluster expansions,
\emph{Lecture Notes in Physics} \textbf{446} (1995).
\href{https://arxiv.org/abs/hep-th/9409094}{arXiv:hep-th/9409094}.
\bibitem{CDP}
W.-K. Chen, N. Dafnis and G. Paouris, Improved H\"older and reverse H\"older inequalities for Gaussian random vectors,
\emph{Adv. Math.} \textbf{280} (2015), Theorem~1.
\href{https://arxiv.org/abs/1306.2410}{arXiv:1306.2410}.

\bibitem{DeSalvo}
S. DeSalvo, Exact sampling algorithms for Latin squares and Sudoku matrices
via probabilistic divide-and-conquer, \href{https://arxiv.org/abs/1502.00235v2}{arXiv:1502.00235v2} (2016).
\bibitem{CM}
E.~R. Canfield and B.~D. McKay, Asymptotic enumeration of dense
$0$--$1$ matrices with equal row sums and equal column sums,
\emph{Electron. J. Combin.} \textbf{12} (2005), R29.
\href{https://doi.org/10.37236/1926}{doi:10.37236/1926}.
\bibitem{TY}
K. Tikhomirov and P. Youssef, The spectral gap of dense random regular graphs,
\emph{Ann. Probab.} \textbf{47} (2019), 362--419;
\href{https://arxiv.org/abs/1610.01765v2}{arXiv:1610.01765v2}.
\bibitem{GKW}
C. Greenhill, J.~H. Kim and N.~C. Wormald, Hamiltonian decompositions of random bipartite regular graphs, proof of Corollary~1.1, equation~(1), \href{https://users.monash.edu.au/~nwormald/papers/hambip.pdf}{author's paper}.
\end{thebibliography}

\end{document}
